\chapter{Оригінальні коментарі до статусу доставки що використувуются в різних компаніях} \label{AppendixA}

\begin{mytable*}{ | p{16cm} | }%
	{Оригінальні коментарі до статусу доставки що використувуются в різних компаніях}%
	{\label{tbl:original_comments}}%
	{\textbf{Коментар до статусу доставки}}
	
	Відмова клієнта прийняти товар \\
	\nopagebreak\quad Відмова клієнта прийняти, замінити чи повернути обладнання \\
	\nopagebreak\quad Не приймає продукцію, експедитор доставив товар на зазначену адресу вчасно \\
	\hline
	Відсутність клієнта \\
	\nopagebreak\quad Клієнт забув про доставку \\
	\nopagebreak\quad Особисті причини (пішов у лікарню, пагане самопочуття, викликали на роботу, пішов на навчання, інше) \\
	\nopagebreak\quad Некоректний номер телефону для звʼязку, недодзвон \\
	\nopagebreak\quad Немає на місці, немає звʼязку \\
	\hline
	Немає грошей \\
	\nopagebreak\quad Відсутність у клієнта коштів, експедитор доставив товар на зазначену адресу вчасно \\
	\hline
	Немає тари \\
	\nopagebreak\quad Відсутність у клієнта тари, експедитор доставив товар на зазначену адресу вчасно \\
	\hline
	Інше \\
	\nopagebreak\quad Експедитор застряг у ліфті \\
	\nopagebreak\quad За станом здоровʼя експедитора \\
	\nopagebreak\quad Заблокований вʼїзд/виїзд автомобіля доставки \\
	\nopagebreak\quad Затримка у попереднього клієнта \\
	\nopagebreak\quad Не працює ліфт, клієнт проживає вище 9-го поверху \\
	\nopagebreak\quad Немає можливості  відмінити замовлення, заявка опрацьована Логістом \\
	\nopagebreak\quad Немає підʼїзду до будинку (будів., дорожні роботи, припаркований транспорт і т.д.) обїхати неможливо \\
	\nopagebreak\quad Непередбачувані обставини (відсутність світла, інтернету тощо) \\
	\nopagebreak\quad Транспортний затор на маршруті \\
	\nopagebreak\quad У клієнта відсутні копії документів \\
	\nopagebreak\quad У клієнта немає печатки \\
	\nopagebreak\quad Транспортний засіб попав у ДТП \\
	\hline
	Погодні умови \\
	\nopagebreak\quad Складні погодні умови (автотранспорт фізично не може доїхати до місця доставки) \\
	\hline
	Невчасна доставка \\
	\nopagebreak\quad Невчасний (запізнення/достроковий) приїзд на точку доставку, клієнт не має змоги прийняти товар \\
	\hline
	Неякісна експедиція \\
	\nopagebreak\quad Відсутність перепускних документів (паспорт, тощо) у експедитора \\
	\nopagebreak\quad Пізно виїхав \\
	\nopagebreak\quad Експедитор не відразу повідомив диспетчера про недовози та їх причини (довідалися в кінці дня) \\
	\nopagebreak\quad Експедитор недотримувався маршруту руху \\
	\nopagebreak\quad Забруднена продукція (пил, бруд, стійкий сторонній запах) \\
	\nopagebreak\quad Завантажений "пересорт" товару \\
	\nopagebreak\quad Не завантажено відповідна до накладних кількість товару на складі \\
	\nopagebreak\quad Немає  накладної \\
	\nopagebreak\quad Недостатня кількість персоналу експедиції \\
	\nopagebreak\quad Пошкодження/втрата бутлів або додаткової продукції, що призначені для клієнта \\
	\hline
	Неякісна логістика \\
	\nopagebreak\quad Багато ранкових замовлень \\
	\nopagebreak\quad Виконання точок іншого маршруту \\
	\nopagebreak\quad Відсутність необхідної кількості транспорту для здійснення доставки \\
	\nopagebreak\quad Відсутність товару на складі \\
	\nopagebreak\quad Затримка завантаження товару на складі за умови вчасної подачі автомобіля \\
	\nopagebreak\quad Неправильне присвоєння сектору/координат \\
	\nopagebreak\quad Неправильно розрахований тонаж та/або обєм продукції (товар не вміщується в авто) \\
	\nopagebreak\quad Похибка при складанні маршруту (адреса точки доставки випадає з логіки загального маршруту) \\
	\hline
	Проблема з транспортом \\
	\nopagebreak\quad Невиїзд запланованого автотранспорту на маршрут \\
	\nopagebreak\quad Поломка ТЗ на маршруті \\
	\hline
	Неправильне замовлення \\
	\nopagebreak\quad Відсутня тара (оператор не попередив), експедитор доставив товар на зазначену в заявці адресу вчасно \\
	\nopagebreak\quad Задвоєне замовлення, половина повертається, експедитор доставив товар на адресу в заявці вчасно \\
	\nopagebreak\quad Клієнт не робив замовлення \\
	\nopagebreak\quad Не вказаний обід, експедитор не може чекати \\
	\nopagebreak\quad Невірно встановлено дату або час, експедитор доставив товар на зазначену в заявці адресу вчасно \\
	\nopagebreak\quad Неправильна адреса, або недостатньо інформації для доставки замовлення \\
	\nopagebreak\quad У документах зазначено товар, який клієнт не замовляв, експедитор доставив згідно накладної вчасно \\
	\hline
	Не встигли вчасно доїхати \\
	\hline
	не знайшли адресу \\
	\hline
	дуже щільний графік \\
	\hline
	запізно виїхали з бази \\
	\hline
	надзвичайні пробки на дорозі \\
	\hline
	надзвичайні погодні умови  \\
	\hline
	дорога до клієнта заблокована \\
	\hline
	технічна несправність автомобіля  \\
	\hline
	аварія автомобіля (ДТП) \\
	\hline
	Клієнт відсутній на місці \\
	\hline
	клієнт не зміг бути вчасно \\
	\hline
	зачинено, телефони не відповідають \\
	\hline
	короткий робочий день, вихідний \\
	\hline
	клієнт зʼїхав з адреси \\
	\hline
	Помилка в замовленні \\
	\nopagebreak\quad не замовляли взагалі \\
	\nopagebreak\quad замовляли на інший день \\
	\nopagebreak\quad замовляли на іншу адресу \\
	\nopagebreak\quad замовляли на інший час \\
	\nopagebreak\quad замовляли інший товар \\
	\hline
	Пошкоджений або відсутній товар \\
	\nopagebreak\quad товар не завантажено  \\
	\nopagebreak\quad товар забруднений \\
	\nopagebreak\quad товар недоукомплектований \\
	\nopagebreak\quad товар не працює (внутрішня поломка) \\
	\nopagebreak\quad товар пошкоджено зовні \\
	\hline
	Скасовано клієнтом (на місці) \\
	\hline
	немає грошей, щоб оплатити замовлення \\
	\hline
	немає потрібної контактної особи \\
	\hline
	закрито доступ до приміщення кліента \\
	\hline
	відсутні бухгалтерські документи \\
	\hline
	помилки в бухгалтерських документах \\
	\hline
	клієнт змінив свої плани щодо замовлення \\
	\hline
	відсутні печатка/пасп.дані клієнта \\
	\hline
	відмовляються сплатити борг (оренду, заставу) \\
	\hline
	не згодні з взаєморозрахунками \\
	\hline
	отримувач не відомий \\
	\hline
	Адреси не існує \\
	\hline
	Невдала спроба доставки \\
	\nopagebreak\quad зміна адреси/переадресація \\
	\nopagebreak\quad за адресою не відчиняють \\
	\nopagebreak\quad не доїхав - нема дорожнього покриття \\
	\nopagebreak\quad не виконав \\
	\nopagebreak\quad з технічних причин \\
	\nopagebreak\quad отримувач переніс дату доставки \\
	\nopagebreak\quad немає грошей \\
	\nopagebreak\quad курʼєр не встигає \\
	\nopagebreak\quad відмова отримувача (утилізація) \\
	\hline
	Заявка на відмову \\
	\nopagebreak\quad брак/пошкодження товару \\
	\nopagebreak\quad невідповідність вкладень замовленню \\
	\nopagebreak\quad нічого не замовляв \\
	\nopagebreak\quad замовлення не актуальне \\
	\nopagebreak\quad немає грошей \\
	\nopagebreak\quad недійсні/некомплектні документи отримувача (кредитна доставка) \\
	\hline
	Відправлення втрачено \\
	\hline
	никого нет дома \\
	\hline
	нет денег \\
	\hline
	говорят, что не заказывали \\
	\hline
	Никого нет \\
	\hline
	Отказ \\
	\hline
	не хотят \\
	\hline
	не просили \\
	\hline
	отобрали и денег не дали \\
	\hline
	Поломка автомобиля \\
	\hline
	Нет доверенности \\
	\hline
	Не успел водитель \\
	\hline
	Нет товароведа \\
	\hline
	Точка закрыта \\
	\hline
	Нет на месте \\
	\hline
	Частичная доставка (бой,рем кулера) \\
	\hline
	Водитель \\
	\nopagebreak\quad Нарушен временной интервал \\
	\nopagebreak\quad Бой балонов \\
	\nopagebreak\quad Не смог найти \\
	\nopagebreak\quad Отмена визита водителем \\
	\hline
	Клиент \\
	\nopagebreak\quad Клиента нет дома \\
	\nopagebreak\quad Отказ клиента на месте \\
	\nopagebreak\quad Отмена визита Клиентом \\
	\nopagebreak\quad Нет тары к возврату \\
	\hline
	Офис \\
	\nopagebreak\quad Неверный адрес \\
	\nopagebreak\quad Отмена визита офисом \\
	\nopagebreak\quad Ошибка офиса  \\
	\hline
	Некорректный маршрут \\
	\hline
	Ошибка склада \\
	\hline
	Форс-мажор \\
	\nopagebreak\quad Поломка авто \\
	\nopagebreak\quad ДТП \\
	\nopagebreak\quad Не выход водителя \\
	\nopagebreak\quad Не поместилось в машину \\
	\nopagebreak\quad Не было места для парковки \\
	\nopagebreak\quad Авто не прошло по габаритам \\
	\nopagebreak\quad Перекрытие дороги \\
	\nopagebreak\quad Нет подъезда к клиенту \\
	\nopagebreak\quad Несвоевременный вызов водителя на маршрут \\
	\nopagebreak\quad Не работает лифт (Выше 5-го этажа) \\
	\hline
	Состояние здоровья водителя \\
	\hline
	Другое \\
	\hline
	Точка закрыта \\
\end{mytable*}%


\chapter{Детальні таблиці точності розпізнавання по реакціях} \label{AppendixB}

\begin{mytable*}{ | c | c | c | c | c | c | }%
	{Результати моделювання конкретних реакцій першого набору даних використовуючи ІРС з послідовностями розміром 1--3}%
	{\label{tbl:detailed_data1_irs13}}%
	{№ Контексту & № Реакції & Точність & Прецизійність & Повнота & F-міра}
	
	1 & 5 & 0.750 & 0.711 & 0.730 & 0.613 \\
\hline
1 & 6 & 0.600 & 0.625 & 0.612 & 0.387 \\
\hline
2 & 6 & 0.138 & 0.917 & 0.240 & 0.094 \\
\hline
2 & 41 & 0.979 & 0.409 & 0.577 & 0.906 \\
\hline
3 & 1 & 0.000 & 0.000 & 0.000 & 0.120 \\
\hline
3 & 32 & 0.000 & 0.000 & 0.000 & 0.097 \\
\hline
3 & 33 & 0.000 & 0.000 & 0.000 & 0.088 \\
\hline
3 & 35 & 0.000 & 0.000 & 0.000 & 0.097 \\
\hline
3 & 55 & 0.000 & 0.000 & 0.000 & 0.101 \\
\hline
3 & 56 & 0.000 & 0.000 & 0.000 & 0.111 \\
\hline
3 & 57 & 0.387 & 1.000 & 0.558 & 0.387 \\
\hline
4 & 8 & 0.500 & 0.083 & 0.143 & 0.231 \\
\hline
4 & 31 & 0.000 & 0.000 & 0.000 & 0.202 \\
\hline
4 & 58 & 0.000 & 0.000 & 0.000 & 0.212 \\
\hline
4 & 59 & 0.370 & 1.000 & 0.540 & 0.356 \\
\hline
5 & 8 & 0.000 & 0.000 & 0.000 & 0.131 \\
\hline
5 & 31 & 0.000 & 0.000 & 0.000 & 0.115 \\
\hline
5 & 60 & 0.000 & 0.000 & 0.000 & 0.115 \\
\hline
5 & 61 & 0.341 & 1.000 & 0.509 & 0.317 \\
\hline
5 & 62 & 0.000 & 0.000 & 0.000 & 0.115 \\
\hline
5 & 63 & 0.308 & 0.105 & 0.157 & 0.208 \\
\hline
6 & 8 & 0.000 & 0.000 & 0.000 & 0.115 \\
\hline
6 & 31 & 0.000 & 0.000 & 0.000 & 0.100 \\
\hline
6 & 64 & 0.396 & 0.456 & 0.424 & 0.378 \\
\hline
6 & 65 & 0.398 & 0.553 & 0.463 & 0.407 \\
\hline
7 & 2 & 0.207 & 0.980 & 0.342 & 0.210 \\
\hline
7 & 32 & 0.000 & 0.000 & 0.000 & 0.090 \\
\hline
7 & 33 & 0.000 & 0.000 & 0.000 & 0.082 \\
\hline
7 & 34 & 0.000 & 0.000 & 0.000 & 0.090 \\
\hline
7 & 35 & 0.000 & 0.000 & 0.000 & 0.090 \\
\hline
7 & 36 & 0.000 & 0.000 & 0.000 & 0.086 \\
\hline
7 & 37 & 0.000 & 0.000 & 0.000 & 0.086 \\
\hline
7 & 38 & 0.000 & 0.000 & 0.000 & 0.090 \\
\hline
7 & 39 & 0.000 & 0.000 & 0.000 & 0.090 \\
\hline
7 & 40 & 0.000 & 0.000 & 0.000 & 0.086 \\
\hline
8 & 8 & 0.000 & 0.000 & 0.000 & 0.111 \\
\hline
8 & 31 & 0.000 & 0.000 & 0.000 & 0.097 \\
\hline
8 & 51 & 0.000 & 0.000 & 0.000 & 0.088 \\
\hline
8 & 52 & 0.000 & 0.000 & 0.000 & 0.184 \\
\hline
8 & 53 & 0.750 & 0.073 & 0.133 & 0.189 \\
\hline
8 & 54 & 0.343 & 1.000 & 0.511 & 0.332 \\
\hline
9 & 8 & 0.375 & 0.125 & 0.188 & 0.183 \\
\hline
9 & 31 & 0.000 & 0.000 & 0.000 & 0.160 \\
\hline
9 & 48 & 0.347 & 1.000 & 0.515 & 0.321 \\
\hline
9 & 49 & 0.000 & 0.000 & 0.000 & 0.176 \\
\hline
9 & 50 & 0.000 & 0.000 & 0.000 & 0.160 \\
\hline
10 & 8 & 0.625 & 0.208 & 0.312 & 0.190 \\
\hline
10 & 31 & 1.000 & 0.048 & 0.091 & 0.167 \\
\hline
10 & 42 & 0.500 & 0.045 & 0.083 & 0.175 \\
\hline
10 & 43 & 0.000 & 0.000 & 0.000 & 0.167 \\
\hline
10 & 44 & 0.322 & 0.974 & 0.484 & 0.302 \\
\hline
11 & 8 & 0.000 & 0.000 & 0.000 & 0.100 \\
\hline
11 & 31 & 0.000 & 0.000 & 0.000 & 0.088 \\
\hline
11 & 45 & 0.299 & 0.562 & 0.391 & 0.372 \\
\hline
11 & 46 & 0.194 & 0.165 & 0.178 & 0.356 \\
\hline
11 & 47 & 0.000 & 0.000 & 0.000 & 0.084 \\
\hline
12 & 3 & 0.182 & 0.462 & 0.261 & 0.129 \\
\hline
12 & 7 & 0.167 & 0.167 & 0.167 & 0.119 \\
\hline
12 & 8 & 0.325 & 0.542 & 0.406 & 0.119 \\
\hline
12 & 28 & 0.257 & 0.360 & 0.300 & 0.124 \\
\hline
12 & 32 & 0.500 & 0.238 & 0.323 & 0.104 \\
\hline
12 & 33 & 0.000 & 0.000 & 0.000 & 0.095 \\
\hline
12 & 34 & 0.600 & 0.286 & 0.387 & 0.104 \\
\hline
12 & 35 & 0.200 & 0.095 & 0.129 & 0.104 \\
\hline
12 & 36 & 0.500 & 0.100 & 0.167 & 0.100 \\
\hline
13 & 4 & 0.000 & 0.000 & 0.000 & 0.275 \\
\hline
13 & 5 & 0.446 & 0.974 & 0.612 & 0.444 \\
\hline
13 & 6 & 1.000 & 0.062 & 0.118 & 0.281 \\
\hline
14 & 8 & 0.296 & 0.667 & 0.410 & 0.173 \\
\hline
14 & 9 & 0.360 & 0.391 & 0.375 & 0.165 \\
\hline
14 & 10 & 0.464 & 0.542 & 0.500 & 0.173 \\
\hline
14 & 11 & 0.769 & 0.417 & 0.541 & 0.173 \\
\hline
14 & 12 & 0.667 & 0.435 & 0.526 & 0.165 \\
\hline
14 & 31 & 0.750 & 0.143 & 0.240 & 0.151 \\
\hline
15 & 8 & 0.500 & 0.042 & 0.077 & 0.151 \\
\hline
15 & 13 & 0.500 & 0.040 & 0.074 & 0.157 \\
\hline
15 & 14 & 0.667 & 0.087 & 0.154 & 0.145 \\
\hline
15 & 15 & 0.284 & 1.000 & 0.442 & 0.264 \\
\hline
15 & 16 & 0.333 & 0.042 & 0.074 & 0.151 \\
\hline
15 & 31 & 1.000 & 0.048 & 0.091 & 0.132 \\
\hline
16 & 8 & 0.395 & 0.708 & 0.507 & 0.209 \\
\hline
16 & 17 & 0.533 & 0.348 & 0.421 & 0.200 \\
\hline
16 & 18 & 0.500 & 0.565 & 0.531 & 0.200 \\
\hline
16 & 19 & 0.519 & 0.583 & 0.549 & 0.209 \\
\hline
16 & 31 & 0.750 & 0.143 & 0.240 & 0.183 \\
\hline
17 & 8 & 0.429 & 0.625 & 0.508 & 0.155 \\
\hline
17 & 20 & 0.333 & 0.348 & 0.340 & 0.148 \\
\hline
17 & 21 & 0.167 & 0.217 & 0.189 & 0.148 \\
\hline
17 & 22 & 0.292 & 0.318 & 0.304 & 0.142 \\
\hline
17 & 23 & 0.087 & 0.095 & 0.091 & 0.135 \\
\hline
17 & 24 & 0.267 & 0.190 & 0.222 & 0.135 \\
\hline
17 & 31 & 0.250 & 0.048 & 0.080 & 0.135 \\
\hline
18 & 8 & 0.315 & 0.708 & 0.436 & 0.218 \\
\hline
18 & 25 & 0.348 & 0.348 & 0.348 & 0.209 \\
\hline
18 & 26 & 0.769 & 0.476 & 0.588 & 0.191 \\
\hline
18 & 27 & 0.643 & 0.429 & 0.514 & 0.191 \\
\hline
18 & 31 & 0.333 & 0.095 & 0.148 & 0.191 \\
\hline
19 & 8 & 0.351 & 0.833 & 0.494 & 0.276 \\
\hline
19 & 29 & 1.000 & 0.381 & 0.552 & 0.241 \\
\hline
19 & 30 & 0.533 & 0.381 & 0.444 & 0.241 \\
\hline
19 & 31 & 0.714 & 0.238 & 0.357 & 0.241 \\
\hline
\specialcell{3cm}{По всій \\ вибірці} & 1 & 0.000 & 0.000 & 0.000 & 0.013 \\
\hline
\specialcell{3cm}{По всій \\ вибірці} & 2 & 0.333 & 0.020 & 0.038 & 0.024 \\
\hline
\specialcell{3cm}{По всій \\ вибірці} & 3 & 0.000 & 0.000 & 0.000 & 0.013 \\
\hline
\specialcell{3cm}{По всій \\ вибірці} & 4 & 0.000 & 0.000 & 0.000 & 0.023 \\
\hline
\specialcell{3cm}{По всій \\ вибірці} & 5 & 0.091 & 0.237 & 0.132 & 0.037 \\
\hline
\specialcell{3cm}{По всій \\ вибірці} & 6 & 1.000 & 0.021 & 0.041 & 0.023 \\
\hline
\specialcell{3cm}{По всій \\ вибірці} & 7 & 0.000 & 0.000 & 0.000 & 0.012 \\
\hline
\specialcell{3cm}{По всій \\ вибірці} & 8 & 0.000 & 0.000 & 0.000 & 0.012 \\
\hline
\specialcell{3cm}{По всій \\ вибірці} & 9 & 0.000 & 0.000 & 0.000 & 0.011 \\
\hline
\specialcell{3cm}{По всій \\ вибірці} & 10 & 0.000 & 0.000 & 0.000 & 0.012 \\
\hline
\specialcell{3cm}{По всій \\ вибірці} & 11 & 0.000 & 0.000 & 0.000 & 0.012 \\
\hline
\specialcell{3cm}{По всій \\ вибірці} & 12 & 0.000 & 0.000 & 0.000 & 0.011 \\
\hline
\specialcell{3cm}{По всій \\ вибірці} & 13 & 0.000 & 0.000 & 0.000 & 0.012 \\
\hline
\specialcell{3cm}{По всій \\ вибірці} & 14 & 0.000 & 0.000 & 0.000 & 0.011 \\
\hline
\specialcell{3cm}{По всій \\ вибірці} & 15 & 0.000 & 0.000 & 0.000 & 0.020 \\
\hline
\specialcell{3cm}{По всій \\ вибірці} & 16 & 0.000 & 0.000 & 0.000 & 0.012 \\
\hline
\specialcell{3cm}{По всій \\ вибірці} & 17 & 0.000 & 0.000 & 0.000 & 0.011 \\
\hline
\specialcell{3cm}{По всій \\ вибірці} & 18 & 0.000 & 0.000 & 0.000 & 0.011 \\
\hline
\specialcell{3cm}{По всій \\ вибірці} & 19 & 0.000 & 0.000 & 0.000 & 0.012 \\
\hline
\specialcell{3cm}{По всій \\ вибірці} & 20 & 0.000 & 0.000 & 0.000 & 0.011 \\
\hline
\specialcell{3cm}{По всій \\ вибірці} & 21 & 0.000 & 0.000 & 0.000 & 0.011 \\
\hline
\specialcell{3cm}{По всій \\ вибірці} & 22 & 0.000 & 0.000 & 0.000 & 0.011 \\
\hline
\specialcell{3cm}{По всій \\ вибірці} & 23 & 0.000 & 0.000 & 0.000 & 0.010 \\
\hline
\specialcell{3cm}{По всій \\ вибірці} & 24 & 0.000 & 0.000 & 0.000 & 0.010 \\
\hline
\specialcell{3cm}{По всій \\ вибірці} & 25 & 0.000 & 0.000 & 0.000 & 0.011 \\
\hline
\specialcell{3cm}{По всій \\ вибірці} & 26 & 0.000 & 0.000 & 0.000 & 0.010 \\
\hline
\specialcell{3cm}{По всій \\ вибірці} & 27 & 0.000 & 0.000 & 0.000 & 0.010 \\
\hline
\specialcell{3cm}{По всій \\ вибірці} & 28 & 0.000 & 0.000 & 0.000 & 0.012 \\
\hline
\specialcell{3cm}{По всій \\ вибірці} & 29 & 0.000 & 0.000 & 0.000 & 0.010 \\
\hline
\specialcell{3cm}{По всій \\ вибірці} & 30 & 0.000 & 0.000 & 0.000 & 0.010 \\
\hline
\specialcell{3cm}{По всій \\ вибірці} & 31 & 0.000 & 0.000 & 0.000 & 0.010 \\
\hline
\specialcell{3cm}{По всій \\ вибірці} & 32 & 0.000 & 0.000 & 0.000 & 0.010 \\
\hline
\specialcell{3cm}{По всій \\ вибірці} & 33 & 0.000 & 0.000 & 0.000 & 0.009 \\
\hline
\specialcell{3cm}{По всій \\ вибірці} & 34 & 0.000 & 0.000 & 0.000 & 0.010 \\
\hline
\specialcell{3cm}{По всій \\ вибірці} & 35 & 0.000 & 0.000 & 0.000 & 0.010 \\
\hline
\specialcell{3cm}{По всій \\ вибірці} & 36 & 0.000 & 0.000 & 0.000 & 0.010 \\
\hline
\specialcell{3cm}{По всій \\ вибірці} & 37 & 0.000 & 0.000 & 0.000 & 0.010 \\
\hline
\specialcell{3cm}{По всій \\ вибірці} & 38 & 0.000 & 0.000 & 0.000 & 0.010 \\
\hline
\specialcell{3cm}{По всій \\ вибірці} & 39 & 0.000 & 0.000 & 0.000 & 0.010 \\
\hline
\specialcell{3cm}{По всій \\ вибірці} & 40 & 0.000 & 0.000 & 0.000 & 0.010 \\
\hline
\specialcell{3cm}{По всій \\ вибірці} & 42 & 0.000 & 0.000 & 0.000 & 0.011 \\
\hline
\specialcell{3cm}{По всій \\ вибірці} & 43 & 0.000 & 0.000 & 0.000 & 0.010 \\
\hline
\specialcell{3cm}{По всій \\ вибірці} & 44 & 0.000 & 0.000 & 0.000 & 0.018 \\
\hline
\specialcell{3cm}{По всій \\ вибірці} & 45 & 0.031 & 0.247 & 0.055 & 0.043 \\
\hline
\specialcell{3cm}{По всій \\ вибірці} & 46 & 0.032 & 0.141 & 0.052 & 0.041 \\
\hline
\specialcell{3cm}{По всій \\ вибірці} & 47 & 0.000 & 0.000 & 0.000 & 0.010 \\
\hline
\specialcell{3cm}{По всій \\ вибірці} & 48 & 0.000 & 0.000 & 0.000 & 0.020 \\
\hline
\specialcell{3cm}{По всій \\ вибірці} & 49 & 0.000 & 0.000 & 0.000 & 0.011 \\
\hline
\specialcell{3cm}{По всій \\ вибірці} & 50 & 0.000 & 0.000 & 0.000 & 0.010 \\
\hline
\specialcell{3cm}{По всій \\ вибірці} & 51 & 0.000 & 0.000 & 0.000 & 0.009 \\
\hline
\specialcell{3cm}{По всій \\ вибірці} & 52 & 0.000 & 0.000 & 0.000 & 0.019 \\
\hline
\specialcell{3cm}{По всій \\ вибірці} & 53 & 0.000 & 0.000 & 0.000 & 0.020 \\
\hline
\specialcell{3cm}{По всій \\ вибірці} & 54 & 0.125 & 0.056 & 0.077 & 0.035 \\
\hline
\specialcell{3cm}{По всій \\ вибірці} & 55 & 0.000 & 0.000 & 0.000 & 0.011 \\
\hline
\specialcell{3cm}{По всій \\ вибірці} & 56 & 0.000 & 0.000 & 0.000 & 0.012 \\
\hline
\specialcell{3cm}{По всій \\ вибірці} & 57 & 0.044 & 0.143 & 0.067 & 0.041 \\
\hline
\specialcell{3cm}{По всій \\ вибірці} & 58 & 0.000 & 0.000 & 0.000 & 0.011 \\
\hline
\specialcell{3cm}{По всій \\ вибірці} & 59 & 0.000 & 0.000 & 0.000 & 0.018 \\
\hline
\specialcell{3cm}{По всій \\ вибірці} & 60 & 0.000 & 0.000 & 0.000 & 0.010 \\
\hline
\specialcell{3cm}{По всій \\ вибірці} & 61 & 0.000 & 0.000 & 0.000 & 0.028 \\
\hline
\specialcell{3cm}{По всій \\ вибірці} & 62 & 0.000 & 0.000 & 0.000 & 0.010 \\
\hline
\specialcell{3cm}{По всій \\ вибірці} & 63 & 0.000 & 0.000 & 0.000 & 0.018 \\
\hline
\specialcell{3cm}{По всій \\ вибірці} & 64 & 0.067 & 0.139 & 0.091 & 0.038 \\
\hline
\specialcell{3cm}{По всій \\ вибірці} & 65 & 0.037 & 0.129 & 0.057 & 0.041 \\
\hline
\specialcell{3cm}{Тестовий \\ контекст} & 1 & 0.000 & 0.000 & 0.000 & 0.257 \\
\hline
\specialcell{3cm}{Тестовий \\ контекст} & 2 & 0.485 & 1.000 & 0.653 & 0.485 \\
\hline
\specialcell{3cm}{Тестовий \\ контекст} & 3 & 0.000 & 0.000 & 0.000 & 0.257 \\
\end{mytable*}

\begin{mytable*}{ | c | c | c | c | c | c | }%
	{Результати моделювання конкретних реакцій першого набору даних використовуючи ІРС з послідовностями розміром 2--4}%
	{\label{tbl:detailed_data1_irs24}}%
	{№ Контексту & № Реакції & Точність & Прецизійність & Повнота & F-міра}
	
	1 & 5 & 0.750 & 0.868 & 0.805 & 0.613 \\
\hline
1 & 6 & 0.759 & 0.458 & 0.571 & 0.387 \\
\hline
2 & 6 & 0.444 & 0.500 & 0.471 & 0.094 \\
\hline
2 & 41 & 0.952 & 0.935 & 0.944 & 0.906 \\
\hline
3 & 1 & 0.000 & 0.000 & 0.000 & 0.120 \\
\hline
3 & 32 & 0.000 & 0.000 & 0.000 & 0.097 \\
\hline
3 & 33 & 0.000 & 0.000 & 0.000 & 0.088 \\
\hline
3 & 35 & 0.000 & 0.000 & 0.000 & 0.097 \\
\hline
3 & 55 & 0.000 & 0.000 & 0.000 & 0.101 \\
\hline
3 & 56 & 0.000 & 0.000 & 0.000 & 0.111 \\
\hline
3 & 57 & 0.384 & 0.988 & 0.553 & 0.387 \\
\hline
4 & 8 & 0.857 & 0.500 & 0.632 & 0.231 \\
\hline
4 & 31 & 0.833 & 0.238 & 0.370 & 0.202 \\
\hline
4 & 58 & 0.800 & 0.182 & 0.296 & 0.212 \\
\hline
4 & 59 & 0.468 & 1.000 & 0.638 & 0.356 \\
\hline
5 & 8 & 0.500 & 0.042 & 0.077 & 0.131 \\
\hline
5 & 31 & 1.000 & 0.143 & 0.250 & 0.115 \\
\hline
5 & 60 & 0.000 & 0.000 & 0.000 & 0.115 \\
\hline
5 & 61 & 0.365 & 1.000 & 0.535 & 0.317 \\
\hline
5 & 62 & 0.000 & 0.000 & 0.000 & 0.115 \\
\hline
5 & 63 & 0.737 & 0.368 & 0.491 & 0.208 \\
\hline
6 & 8 & 0.000 & 0.000 & 0.000 & 0.115 \\
\hline
6 & 31 & 0.000 & 0.000 & 0.000 & 0.100 \\
\hline
6 & 64 & 0.551 & 0.684 & 0.610 & 0.378 \\
\hline
6 & 65 & 0.577 & 0.753 & 0.653 & 0.407 \\
\hline
7 & 2 & 0.216 & 0.980 & 0.354 & 0.210 \\
\hline
7 & 32 & 0.667 & 0.095 & 0.167 & 0.090 \\
\hline
7 & 33 & 0.000 & 0.000 & 0.000 & 0.082 \\
\hline
7 & 34 & 0.000 & 0.000 & 0.000 & 0.090 \\
\hline
7 & 35 & 1.000 & 0.048 & 0.091 & 0.090 \\
\hline
7 & 36 & 0.000 & 0.000 & 0.000 & 0.086 \\
\hline
7 & 37 & 1.000 & 0.100 & 0.182 & 0.086 \\
\hline
7 & 38 & 0.000 & 0.000 & 0.000 & 0.090 \\
\hline
7 & 39 & 0.000 & 0.000 & 0.000 & 0.090 \\
\hline
7 & 40 & 1.000 & 0.050 & 0.095 & 0.086 \\
\hline
8 & 8 & 0.000 & 0.000 & 0.000 & 0.111 \\
\hline
8 & 31 & 1.000 & 0.048 & 0.091 & 0.097 \\
\hline
8 & 51 & 0.000 & 0.000 & 0.000 & 0.088 \\
\hline
8 & 52 & 0.500 & 0.050 & 0.091 & 0.184 \\
\hline
8 & 53 & 0.600 & 0.146 & 0.235 & 0.189 \\
\hline
8 & 54 & 0.358 & 1.000 & 0.527 & 0.332 \\
\hline
9 & 8 & 0.722 & 0.542 & 0.619 & 0.183 \\
\hline
9 & 31 & 1.000 & 0.190 & 0.320 & 0.160 \\
\hline
9 & 48 & 0.438 & 1.000 & 0.609 & 0.321 \\
\hline
9 & 49 & 0.857 & 0.261 & 0.400 & 0.176 \\
\hline
9 & 50 & 0.667 & 0.190 & 0.296 & 0.160 \\
\hline
10 & 8 & 0.650 & 0.542 & 0.591 & 0.190 \\
\hline
10 & 31 & 1.000 & 0.190 & 0.320 & 0.167 \\
\hline
10 & 42 & 0.800 & 0.182 & 0.296 & 0.175 \\
\hline
10 & 43 & 1.000 & 0.476 & 0.645 & 0.167 \\
\hline
10 & 44 & 0.414 & 0.947 & 0.576 & 0.302 \\
\hline
11 & 8 & 0.000 & 0.000 & 0.000 & 0.100 \\
\hline
11 & 31 & 0.000 & 0.000 & 0.000 & 0.088 \\
\hline
11 & 45 & 0.496 & 0.787 & 0.609 & 0.372 \\
\hline
11 & 46 & 0.541 & 0.624 & 0.579 & 0.356 \\
\hline
11 & 47 & 0.000 & 0.000 & 0.000 & 0.084 \\
\hline
12 & 3 & 0.412 & 0.808 & 0.545 & 0.129 \\
\hline
12 & 7 & 0.235 & 0.167 & 0.195 & 0.119 \\
\hline
12 & 8 & 0.514 & 0.750 & 0.610 & 0.119 \\
\hline
12 & 28 & 0.588 & 0.800 & 0.678 & 0.124 \\
\hline
12 & 32 & 0.500 & 0.381 & 0.432 & 0.104 \\
\hline
12 & 33 & 0.500 & 0.105 & 0.174 & 0.095 \\
\hline
12 & 34 & 0.400 & 0.381 & 0.390 & 0.104 \\
\hline
12 & 35 & 0.667 & 0.476 & 0.556 & 0.104 \\
\hline
12 & 36 & 0.778 & 0.350 & 0.483 & 0.100 \\
\hline
13 & 4 & 0.533 & 0.170 & 0.258 & 0.275 \\
\hline
13 & 5 & 0.482 & 0.895 & 0.627 & 0.444 \\
\hline
13 & 6 & 0.778 & 0.146 & 0.246 & 0.281 \\
\hline
14 & 8 & 0.439 & 0.750 & 0.554 & 0.173 \\
\hline
14 & 9 & 0.550 & 0.478 & 0.512 & 0.165 \\
\hline
14 & 10 & 0.588 & 0.833 & 0.690 & 0.173 \\
\hline
14 & 11 & 0.947 & 0.750 & 0.837 & 0.173 \\
\hline
14 & 12 & 0.833 & 0.652 & 0.732 & 0.165 \\
\hline
14 & 31 & 0.857 & 0.286 & 0.429 & 0.151 \\
\hline
15 & 8 & 0.167 & 0.083 & 0.111 & 0.151 \\
\hline
15 & 13 & 0.778 & 0.280 & 0.412 & 0.157 \\
\hline
15 & 14 & 0.500 & 0.348 & 0.410 & 0.145 \\
\hline
15 & 15 & 0.418 & 0.976 & 0.586 & 0.264 \\
\hline
15 & 16 & 0.500 & 0.375 & 0.429 & 0.151 \\
\hline
15 & 31 & 1.000 & 0.238 & 0.385 & 0.132 \\
\hline
16 & 8 & 0.513 & 0.833 & 0.635 & 0.209 \\
\hline
16 & 17 & 0.833 & 0.217 & 0.345 & 0.200 \\
\hline
16 & 18 & 0.667 & 0.696 & 0.681 & 0.200 \\
\hline
16 & 19 & 0.514 & 0.792 & 0.623 & 0.209 \\
\hline
16 & 31 & 0.778 & 0.333 & 0.467 & 0.183 \\
\hline
17 & 8 & 0.594 & 0.792 & 0.679 & 0.155 \\
\hline
17 & 20 & 0.458 & 0.478 & 0.468 & 0.148 \\
\hline
17 & 21 & 0.281 & 0.391 & 0.327 & 0.148 \\
\hline
17 & 22 & 0.214 & 0.273 & 0.240 & 0.142 \\
\hline
17 & 23 & 0.095 & 0.095 & 0.095 & 0.135 \\
\hline
17 & 24 & 0.667 & 0.476 & 0.556 & 0.135 \\
\hline
17 & 31 & 1.000 & 0.143 & 0.250 & 0.135 \\
\hline
18 & 8 & 0.388 & 0.792 & 0.521 & 0.218 \\
\hline
18 & 25 & 0.684 & 0.565 & 0.619 & 0.209 \\
\hline
18 & 26 & 0.750 & 0.714 & 0.732 & 0.191 \\
\hline
18 & 27 & 0.786 & 0.524 & 0.629 & 0.191 \\
\hline
18 & 31 & 0.750 & 0.286 & 0.414 & 0.191 \\
\hline
19 & 8 & 0.543 & 0.792 & 0.644 & 0.276 \\
\hline
19 & 29 & 0.941 & 0.762 & 0.842 & 0.241 \\
\hline
19 & 30 & 0.571 & 0.762 & 0.653 & 0.241 \\
\hline
19 & 31 & 0.857 & 0.286 & 0.429 & 0.241 \\
\hline
\specialcell{3cm}{По всій \\ вибірці} & 1 & 0.000 & 0.000 & 0.000 & 0.013 \\
\hline
\specialcell{3cm}{По всій \\ вибірці} & 2 & 0.600 & 0.061 & 0.111 & 0.024 \\
\hline
\specialcell{3cm}{По всій \\ вибірці} & 3 & 0.000 & 0.000 & 0.000 & 0.013 \\
\hline
\specialcell{3cm}{По всій \\ вибірці} & 4 & 0.000 & 0.000 & 0.000 & 0.023 \\
\hline
\specialcell{3cm}{По всій \\ вибірці} & 5 & 0.233 & 0.355 & 0.281 & 0.037 \\
\hline
\specialcell{3cm}{По всій \\ вибірці} & 6 & 1.000 & 0.021 & 0.041 & 0.023 \\
\hline
\specialcell{3cm}{По всій \\ вибірці} & 7 & 0.000 & 0.000 & 0.000 & 0.012 \\
\hline
\specialcell{3cm}{По всій \\ вибірці} & 8 & 0.000 & 0.000 & 0.000 & 0.012 \\
\hline
\specialcell{3cm}{По всій \\ вибірці} & 9 & 0.000 & 0.000 & 0.000 & 0.011 \\
\hline
\specialcell{3cm}{По всій \\ вибірці} & 10 & 0.000 & 0.000 & 0.000 & 0.012 \\
\hline
\specialcell{3cm}{По всій \\ вибірці} & 11 & 0.000 & 0.000 & 0.000 & 0.012 \\
\hline
\specialcell{3cm}{По всій \\ вибірці} & 12 & 0.000 & 0.000 & 0.000 & 0.011 \\
\hline
\specialcell{3cm}{По всій \\ вибірці} & 13 & 0.000 & 0.000 & 0.000 & 0.012 \\
\hline
\specialcell{3cm}{По всій \\ вибірці} & 14 & 0.000 & 0.000 & 0.000 & 0.011 \\
\hline
\specialcell{3cm}{По всій \\ вибірці} & 15 & 0.000 & 0.000 & 0.000 & 0.020 \\
\hline
\specialcell{3cm}{По всій \\ вибірці} & 16 & 0.000 & 0.000 & 0.000 & 0.012 \\
\hline
\specialcell{3cm}{По всій \\ вибірці} & 17 & 0.000 & 0.000 & 0.000 & 0.011 \\
\hline
\specialcell{3cm}{По всій \\ вибірці} & 18 & 0.000 & 0.000 & 0.000 & 0.011 \\
\hline
\specialcell{3cm}{По всій \\ вибірці} & 19 & 0.000 & 0.000 & 0.000 & 0.012 \\
\hline
\specialcell{3cm}{По всій \\ вибірці} & 20 & 0.000 & 0.000 & 0.000 & 0.011 \\
\hline
\specialcell{3cm}{По всій \\ вибірці} & 21 & 0.000 & 0.000 & 0.000 & 0.011 \\
\hline
\specialcell{3cm}{По всій \\ вибірці} & 22 & 0.000 & 0.000 & 0.000 & 0.011 \\
\hline
\specialcell{3cm}{По всій \\ вибірці} & 23 & 0.000 & 0.000 & 0.000 & 0.010 \\
\hline
\specialcell{3cm}{По всій \\ вибірці} & 24 & 0.000 & 0.000 & 0.000 & 0.010 \\
\hline
\specialcell{3cm}{По всій \\ вибірці} & 25 & 0.000 & 0.000 & 0.000 & 0.011 \\
\hline
\specialcell{3cm}{По всій \\ вибірці} & 26 & 0.000 & 0.000 & 0.000 & 0.010 \\
\hline
\specialcell{3cm}{По всій \\ вибірці} & 27 & 0.000 & 0.000 & 0.000 & 0.010 \\
\hline
\specialcell{3cm}{По всій \\ вибірці} & 28 & 0.000 & 0.000 & 0.000 & 0.012 \\
\hline
\specialcell{3cm}{По всій \\ вибірці} & 29 & 0.000 & 0.000 & 0.000 & 0.010 \\
\hline
\specialcell{3cm}{По всій \\ вибірці} & 30 & 0.000 & 0.000 & 0.000 & 0.010 \\
\hline
\specialcell{3cm}{По всій \\ вибірці} & 31 & 0.000 & 0.000 & 0.000 & 0.010 \\
\hline
\specialcell{3cm}{По всій \\ вибірці} & 32 & 0.000 & 0.000 & 0.000 & 0.010 \\
\hline
\specialcell{3cm}{По всій \\ вибірці} & 33 & 0.000 & 0.000 & 0.000 & 0.009 \\
\hline
\specialcell{3cm}{По всій \\ вибірці} & 34 & 0.000 & 0.000 & 0.000 & 0.010 \\
\hline
\specialcell{3cm}{По всій \\ вибірці} & 35 & 0.000 & 0.000 & 0.000 & 0.010 \\
\hline
\specialcell{3cm}{По всій \\ вибірці} & 36 & 0.000 & 0.000 & 0.000 & 0.010 \\
\hline
\specialcell{3cm}{По всій \\ вибірці} & 37 & 0.000 & 0.000 & 0.000 & 0.010 \\
\hline
\specialcell{3cm}{По всій \\ вибірці} & 38 & 0.000 & 0.000 & 0.000 & 0.010 \\
\hline
\specialcell{3cm}{По всій \\ вибірці} & 39 & 0.000 & 0.000 & 0.000 & 0.010 \\
\hline
\specialcell{3cm}{По всій \\ вибірці} & 40 & 0.000 & 0.000 & 0.000 & 0.010 \\
\hline
\specialcell{3cm}{По всій \\ вибірці} & 42 & 0.000 & 0.000 & 0.000 & 0.011 \\
\hline
\specialcell{3cm}{По всій \\ вибірці} & 43 & 0.000 & 0.000 & 0.000 & 0.010 \\
\hline
\specialcell{3cm}{По всій \\ вибірці} & 44 & 0.000 & 0.000 & 0.000 & 0.018 \\
\hline
\specialcell{3cm}{По всій \\ вибірці} & 45 & 0.083 & 0.618 & 0.146 & 0.043 \\
\hline
\specialcell{3cm}{По всій \\ вибірці} & 46 & 0.104 & 0.553 & 0.175 & 0.041 \\
\hline
\specialcell{3cm}{По всій \\ вибірці} & 47 & 0.000 & 0.000 & 0.000 & 0.010 \\
\hline
\specialcell{3cm}{По всій \\ вибірці} & 48 & 0.000 & 0.000 & 0.000 & 0.020 \\
\hline
\specialcell{3cm}{По всій \\ вибірці} & 49 & 0.000 & 0.000 & 0.000 & 0.011 \\
\hline
\specialcell{3cm}{По всій \\ вибірці} & 50 & 0.000 & 0.000 & 0.000 & 0.010 \\
\hline
\specialcell{3cm}{По всій \\ вибірці} & 51 & 0.000 & 0.000 & 0.000 & 0.009 \\
\hline
\specialcell{3cm}{По всій \\ вибірці} & 52 & 0.000 & 0.000 & 0.000 & 0.019 \\
\hline
\specialcell{3cm}{По всій \\ вибірці} & 53 & 1.000 & 0.024 & 0.048 & 0.020 \\
\hline
\specialcell{3cm}{По всій \\ вибірці} & 54 & 0.228 & 0.319 & 0.266 & 0.035 \\
\hline
\specialcell{3cm}{По всій \\ вибірці} & 55 & 0.000 & 0.000 & 0.000 & 0.011 \\
\hline
\specialcell{3cm}{По всій \\ вибірці} & 56 & 0.000 & 0.000 & 0.000 & 0.012 \\
\hline
\specialcell{3cm}{По всій \\ вибірці} & 57 & 0.131 & 0.429 & 0.201 & 0.041 \\
\hline
\specialcell{3cm}{По всій \\ вибірці} & 58 & 0.000 & 0.000 & 0.000 & 0.011 \\
\hline
\specialcell{3cm}{По всій \\ вибірці} & 59 & 0.000 & 0.000 & 0.000 & 0.018 \\
\hline
\specialcell{3cm}{По всій \\ вибірці} & 60 & 0.000 & 0.000 & 0.000 & 0.010 \\
\hline
\specialcell{3cm}{По всій \\ вибірці} & 61 & 0.375 & 0.103 & 0.162 & 0.028 \\
\hline
\specialcell{3cm}{По всій \\ вибірці} & 62 & 0.000 & 0.000 & 0.000 & 0.010 \\
\hline
\specialcell{3cm}{По всій \\ вибірці} & 63 & 0.000 & 0.000 & 0.000 & 0.018 \\
\hline
\specialcell{3cm}{По всій \\ вибірці} & 64 & 0.114 & 0.228 & 0.152 & 0.038 \\
\hline
\specialcell{3cm}{По всій \\ вибірці} & 65 & 0.126 & 0.412 & 0.193 & 0.041 \\
\hline
\specialcell{3cm}{Тестовий \\ контекст} & 1 & 0.000 & 0.000 & 0.000 & 0.257 \\
\hline
\specialcell{3cm}{Тестовий \\ контекст} & 2 & 0.480 & 0.980 & 0.644 & 0.485 \\
\hline
\specialcell{3cm}{Тестовий \\ контекст} & 3 & 0.000 & 0.000 & 0.000 & 0.257 \\
\end{mytable*}

\begin{mytable*}{ | c | c | c | c | c | c | }%
	{Результати моделювання конкретних реакцій першого набору даних використовуючи ЗНМ з послідовностями розміром 2--4}%
	{\label{tbl:detailed_data1_cnn}}%
	{№ Контексту & № Реакції & Точність & Прецизійність & Повнота & F-міра}
	
	1 & 5 & 0.877 & 0.934 & 0.904 & 0.613 \\
\hline
1 & 6 & 0.884 & 0.792 & 0.835 & 0.387 \\
\hline
2 & 6 & 0.906 & 0.604 & 0.725 & 0.094 \\
\hline
2 & 41 & 0.960 & 0.994 & 0.977 & 0.906 \\
\hline
3 & 1 & 0.821 & 0.885 & 0.852 & 0.120 \\
\hline
3 & 32 & 1.000 & 0.714 & 0.833 & 0.097 \\
\hline
3 & 33 & 0.842 & 0.842 & 0.842 & 0.088 \\
\hline
3 & 35 & 0.941 & 0.762 & 0.842 & 0.097 \\
\hline
3 & 55 & 0.938 & 0.682 & 0.789 & 0.101 \\
\hline
3 & 56 & 0.789 & 0.625 & 0.698 & 0.111 \\
\hline
3 & 57 & 0.757 & 0.929 & 0.834 & 0.387 \\
\hline
4 & 8 & 0.750 & 0.750 & 0.750 & 0.231 \\
\hline
4 & 31 & 0.933 & 0.667 & 0.778 & 0.202 \\
\hline
4 & 58 & 0.913 & 0.955 & 0.933 & 0.212 \\
\hline
4 & 59 & 0.833 & 0.946 & 0.886 & 0.356 \\
\hline
5 & 8 & 0.826 & 0.792 & 0.809 & 0.131 \\
\hline
5 & 31 & 0.857 & 0.571 & 0.686 & 0.115 \\
\hline
5 & 60 & 0.700 & 0.333 & 0.452 & 0.115 \\
\hline
5 & 61 & 0.753 & 0.948 & 0.840 & 0.317 \\
\hline
5 & 62 & 0.944 & 0.810 & 0.872 & 0.115 \\
\hline
5 & 63 & 0.800 & 0.947 & 0.867 & 0.208 \\
\hline
6 & 8 & 0.905 & 0.792 & 0.844 & 0.115 \\
\hline
6 & 31 & 0.800 & 0.571 & 0.667 & 0.100 \\
\hline
6 & 64 & 0.864 & 0.886 & 0.875 & 0.378 \\
\hline
6 & 65 & 0.837 & 0.906 & 0.870 & 0.407 \\
\hline
7 & 2 & 0.741 & 0.878 & 0.804 & 0.210 \\
\hline
7 & 32 & 0.889 & 0.762 & 0.821 & 0.090 \\
\hline
7 & 33 & 0.938 & 0.789 & 0.857 & 0.082 \\
\hline
7 & 34 & 0.789 & 0.714 & 0.750 & 0.090 \\
\hline
7 & 35 & 0.895 & 0.810 & 0.850 & 0.090 \\
\hline
7 & 36 & 0.889 & 0.800 & 0.842 & 0.086 \\
\hline
7 & 37 & 0.882 & 0.750 & 0.811 & 0.086 \\
\hline
7 & 38 & 0.765 & 0.619 & 0.684 & 0.090 \\
\hline
7 & 39 & 0.680 & 0.810 & 0.739 & 0.090 \\
\hline
7 & 40 & 0.615 & 0.800 & 0.696 & 0.086 \\
\hline
8 & 8 & 0.913 & 0.875 & 0.894 & 0.111 \\
\hline
8 & 31 & 0.875 & 0.667 & 0.757 & 0.097 \\
\hline
8 & 51 & 1.000 & 0.947 & 0.973 & 0.088 \\
\hline
8 & 52 & 0.868 & 0.825 & 0.846 & 0.184 \\
\hline
8 & 53 & 0.889 & 0.780 & 0.831 & 0.189 \\
\hline
8 & 54 & 0.826 & 0.986 & 0.899 & 0.332 \\
\hline
9 & 8 & 0.808 & 0.875 & 0.840 & 0.183 \\
\hline
9 & 31 & 0.933 & 0.667 & 0.778 & 0.160 \\
\hline
9 & 48 & 0.857 & 1.000 & 0.923 & 0.321 \\
\hline
9 & 49 & 0.857 & 0.783 & 0.818 & 0.176 \\
\hline
9 & 50 & 0.900 & 0.857 & 0.878 & 0.160 \\
\hline
10 & 8 & 0.808 & 0.875 & 0.840 & 0.190 \\
\hline
10 & 31 & 0.923 & 0.571 & 0.706 & 0.167 \\
\hline
10 & 42 & 0.850 & 0.773 & 0.810 & 0.175 \\
\hline
10 & 43 & 0.840 & 1.000 & 0.913 & 0.167 \\
\hline
10 & 44 & 0.786 & 0.868 & 0.825 & 0.302 \\
\hline
11 & 8 & 0.875 & 0.875 & 0.875 & 0.100 \\
\hline
11 & 31 & 0.929 & 0.619 & 0.743 & 0.088 \\
\hline
11 & 45 & 0.904 & 0.955 & 0.929 & 0.372 \\
\hline
11 & 46 & 0.900 & 0.953 & 0.926 & 0.356 \\
\hline
11 & 47 & 1.000 & 0.850 & 0.919 & 0.084 \\
\hline
12 & 3 & 0.833 & 0.769 & 0.800 & 0.129 \\
\hline
12 & 7 & 0.783 & 0.750 & 0.766 & 0.119 \\
\hline
12 & 8 & 0.840 & 0.875 & 0.857 & 0.119 \\
\hline
12 & 28 & 0.885 & 0.920 & 0.902 & 0.124 \\
\hline
12 & 32 & 0.947 & 0.857 & 0.900 & 0.104 \\
\hline
12 & 33 & 0.826 & 1.000 & 0.905 & 0.095 \\
\hline
12 & 34 & 0.720 & 0.857 & 0.783 & 0.104 \\
\hline
12 & 35 & 0.895 & 0.810 & 0.850 & 0.104 \\
\hline
12 & 36 & 0.941 & 0.800 & 0.865 & 0.100 \\
\hline
13 & 4 & 0.911 & 0.872 & 0.891 & 0.275 \\
\hline
13 & 5 & 0.807 & 0.882 & 0.843 & 0.444 \\
\hline
13 & 6 & 0.860 & 0.771 & 0.813 & 0.281 \\
\hline
14 & 8 & 0.679 & 0.792 & 0.731 & 0.173 \\
\hline
14 & 9 & 0.842 & 0.696 & 0.762 & 0.165 \\
\hline
14 & 10 & 0.920 & 0.958 & 0.939 & 0.173 \\
\hline
14 & 11 & 0.885 & 0.958 & 0.920 & 0.173 \\
\hline
14 & 12 & 0.846 & 0.957 & 0.898 & 0.165 \\
\hline
14 & 31 & 0.933 & 0.667 & 0.778 & 0.151 \\
\hline
15 & 8 & 0.667 & 0.667 & 0.667 & 0.151 \\
\hline
15 & 13 & 0.923 & 0.960 & 0.941 & 0.157 \\
\hline
15 & 14 & 0.773 & 0.739 & 0.756 & 0.145 \\
\hline
15 & 15 & 0.891 & 0.976 & 0.932 & 0.264 \\
\hline
15 & 16 & 0.917 & 0.917 & 0.917 & 0.151 \\
\hline
15 & 31 & 0.882 & 0.714 & 0.789 & 0.132 \\
\hline
16 & 8 & 0.870 & 0.833 & 0.851 & 0.209 \\
\hline
16 & 17 & 0.842 & 0.696 & 0.762 & 0.200 \\
\hline
16 & 18 & 0.769 & 0.870 & 0.816 & 0.200 \\
\hline
16 & 19 & 0.767 & 0.958 & 0.852 & 0.209 \\
\hline
16 & 31 & 0.882 & 0.714 & 0.789 & 0.183 \\
\hline
17 & 8 & 0.741 & 0.833 & 0.784 & 0.155 \\
\hline
17 & 20 & 0.741 & 0.870 & 0.800 & 0.148 \\
\hline
17 & 21 & 0.667 & 0.435 & 0.526 & 0.148 \\
\hline
17 & 22 & 0.556 & 0.682 & 0.612 & 0.142 \\
\hline
17 & 23 & 0.750 & 0.714 & 0.732 & 0.135 \\
\hline
17 & 24 & 0.947 & 0.857 & 0.900 & 0.135 \\
\hline
17 & 31 & 0.700 & 0.667 & 0.683 & 0.135 \\
\hline
18 & 8 & 0.741 & 0.833 & 0.784 & 0.218 \\
\hline
18 & 25 & 0.905 & 0.826 & 0.864 & 0.209 \\
\hline
18 & 26 & 0.900 & 0.857 & 0.878 & 0.191 \\
\hline
18 & 27 & 0.840 & 1.000 & 0.913 & 0.191 \\
\hline
18 & 31 & 0.824 & 0.667 & 0.737 & 0.191 \\
\hline
19 & 8 & 0.833 & 0.833 & 0.833 & 0.276 \\
\hline
19 & 29 & 0.913 & 1.000 & 0.955 & 0.241 \\
\hline
19 & 30 & 0.808 & 1.000 & 0.894 & 0.241 \\
\hline
19 & 31 & 0.929 & 0.619 & 0.743 & 0.241 \\
\hline
\specialcell{3cm}{По всій \\ вибірці} & 1 & 0.655 & 0.776 & 0.710 & 0.024 \\
\hline
\specialcell{3cm}{По всій \\ вибірці} & 2 & 0.895 & 0.654 & 0.756 & 0.013 \\
\hline
\specialcell{3cm}{По всій \\ вибірці} & 3 & 0.653 & 0.681 & 0.667 & 0.023 \\
\hline
\specialcell{3cm}{По всій \\ вибірці} & 4 & 0.523 & 0.605 & 0.561 & 0.037 \\
\hline
\specialcell{3cm}{По всій \\ вибірці} & 5 & 0.652 & 0.625 & 0.638 & 0.023 \\
\hline
\specialcell{3cm}{По всій \\ вибірці} & 6 & 0.737 & 0.583 & 0.651 & 0.012 \\
\hline
\specialcell{3cm}{По всій \\ вибірці} & 7 & 0.625 & 0.625 & 0.625 & 0.012 \\
\hline
\specialcell{3cm}{По всій \\ вибірці} & 8 & 0.632 & 0.522 & 0.571 & 0.011 \\
\hline
\specialcell{3cm}{По всій \\ вибірці} & 9 & 0.731 & 0.792 & 0.760 & 0.012 \\
\hline
\specialcell{3cm}{По всій \\ вибірці} & 10 & 0.773 & 0.708 & 0.739 & 0.012 \\
\hline
\specialcell{3cm}{По всій \\ вибірці} & 11 & 0.636 & 0.609 & 0.622 & 0.011 \\
\hline
\specialcell{3cm}{По всій \\ вибірці} & 12 & 0.810 & 0.680 & 0.739 & 0.012 \\
\hline
\specialcell{3cm}{По всій \\ вибірці} & 13 & 0.706 & 0.522 & 0.600 & 0.011 \\
\hline
\specialcell{3cm}{По всій \\ вибірці} & 14 & 0.509 & 0.667 & 0.577 & 0.020 \\
\hline
\specialcell{3cm}{По всій \\ вибірці} & 15 & 0.643 & 0.750 & 0.692 & 0.012 \\
\hline
\specialcell{3cm}{По всій \\ вибірці} & 16 & 0.400 & 0.261 & 0.316 & 0.011 \\
\hline
\specialcell{3cm}{По всій \\ вибірці} & 17 & 0.700 & 0.609 & 0.651 & 0.011 \\
\hline
\specialcell{3cm}{По всій \\ вибірці} & 18 & 0.594 & 0.792 & 0.679 & 0.012 \\
\hline
\specialcell{3cm}{По всій \\ вибірці} & 19 & 0.727 & 0.696 & 0.711 & 0.011 \\
\hline
\specialcell{3cm}{По всій \\ вибірці} & 20 & 0.714 & 0.435 & 0.541 & 0.011 \\
\hline
\specialcell{3cm}{По всій \\ вибірці} & 21 & 0.278 & 0.227 & 0.250 & 0.011 \\
\hline
\specialcell{3cm}{По всій \\ вибірці} & 22 & 0.722 & 0.619 & 0.667 & 0.010 \\
\hline
\specialcell{3cm}{По всій \\ вибірці} & 23 & 0.591 & 0.619 & 0.605 & 0.010 \\
\hline
\specialcell{3cm}{По всій \\ вибірці} & 24 & 0.421 & 0.348 & 0.381 & 0.011 \\
\hline
\specialcell{3cm}{По всій \\ вибірці} & 25 & 0.824 & 0.667 & 0.737 & 0.010 \\
\hline
\specialcell{3cm}{По всій \\ вибірці} & 26 & 0.500 & 0.381 & 0.432 & 0.010 \\
\hline
\specialcell{3cm}{По всій \\ вибірці} & 27 & 0.724 & 0.840 & 0.778 & 0.012 \\
\hline
\specialcell{3cm}{По всій \\ вибірці} & 28 & 0.760 & 0.905 & 0.826 & 0.010 \\
\hline
\specialcell{3cm}{По всій \\ вибірці} & 29 & 0.850 & 0.810 & 0.829 & 0.010 \\
\hline
\specialcell{3cm}{По всій \\ вибірці} & 30 & 0.455 & 0.238 & 0.312 & 0.010 \\
\hline
\specialcell{3cm}{По всій \\ вибірці} & 31 & 0.800 & 0.762 & 0.780 & 0.010 \\
\hline
\specialcell{3cm}{По всій \\ вибірці} & 32 & 0.812 & 0.684 & 0.743 & 0.009 \\
\hline
\specialcell{3cm}{По всій \\ вибірці} & 33 & 0.588 & 0.476 & 0.526 & 0.010 \\
\hline
\specialcell{3cm}{По всій \\ вибірці} & 34 & 0.733 & 0.524 & 0.611 & 0.010 \\
\hline
\specialcell{3cm}{По всій \\ вибірці} & 35 & 0.750 & 0.450 & 0.563 & 0.010 \\
\hline
\specialcell{3cm}{По всій \\ вибірці} & 36 & 0.875 & 0.700 & 0.778 & 0.010 \\
\hline
\specialcell{3cm}{По всій \\ вибірці} & 37 & 0.692 & 0.429 & 0.529 & 0.010 \\
\hline
\specialcell{3cm}{По всій \\ вибірці} & 38 & 0.625 & 0.714 & 0.667 & 0.010 \\
\hline
\specialcell{3cm}{По всій \\ вибірці} & 39 & 0.565 & 0.650 & 0.605 & 0.010 \\
\hline
\specialcell{3cm}{По всій \\ вибірці} & 40 & 0.889 & 0.364 & 0.516 & 0.011 \\
\hline
\specialcell{3cm}{По всій \\ вибірці} & 41 & 0.875 & 0.667 & 0.757 & 0.010 \\
\hline
\specialcell{3cm}{По всій \\ вибірці} & 42 & 0.750 & 0.553 & 0.636 & 0.018 \\
\hline
\specialcell{3cm}{По всій \\ вибірці} & 43 & 0.608 & 0.854 & 0.710 & 0.043 \\
\hline
\specialcell{3cm}{По всій \\ вибірці} & 44 & 0.673 & 0.776 & 0.721 & 0.041 \\
\hline
\specialcell{3cm}{По всій \\ вибірці} & 45 & 0.778 & 0.700 & 0.737 & 0.010 \\
\hline
\specialcell{3cm}{По всій \\ вибірці} & 46 & 0.683 & 0.667 & 0.675 & 0.020 \\
\hline
\specialcell{3cm}{По всій \\ вибірці} & 47 & 0.944 & 0.739 & 0.829 & 0.011 \\
\hline
\specialcell{3cm}{По всій \\ вибірці} & 48 & 0.500 & 0.524 & 0.512 & 0.010 \\
\hline
\specialcell{3cm}{По всій \\ вибірці} & 49 & 0.714 & 0.789 & 0.750 & 0.009 \\
\hline
\specialcell{3cm}{По всій \\ вибірці} & 50 & 0.579 & 0.550 & 0.564 & 0.019 \\
\hline
\specialcell{3cm}{По всій \\ вибірці} & 51 & 0.767 & 0.561 & 0.648 & 0.020 \\
\hline
\specialcell{3cm}{По всій \\ вибірці} & 52 & 0.663 & 0.792 & 0.722 & 0.035 \\
\hline
\specialcell{3cm}{По всій \\ вибірці} & 53 & 0.733 & 0.500 & 0.595 & 0.011 \\
\hline
\specialcell{3cm}{По всій \\ вибірці} & 54 & 0.789 & 0.625 & 0.698 & 0.012 \\
\hline
\specialcell{3cm}{По всій \\ вибірці} & 55 & 0.529 & 0.655 & 0.585 & 0.041 \\
\hline
\specialcell{3cm}{По всій \\ вибірці} & 56 & 0.765 & 0.591 & 0.667 & 0.011 \\
\hline
\specialcell{3cm}{По всій \\ вибірці} & 57 & 0.844 & 0.730 & 0.783 & 0.018 \\
\hline
\specialcell{3cm}{По всій \\ вибірці} & 58 & 0.583 & 0.333 & 0.424 & 0.010 \\
\hline
\specialcell{3cm}{По всій \\ вибірці} & 59 & 0.691 & 0.810 & 0.746 & 0.028 \\
\hline
\specialcell{3cm}{По всій \\ вибірці} & 60 & 0.786 & 0.524 & 0.629 & 0.010 \\
\hline
\specialcell{3cm}{По всій \\ вибірці} & 61 & 0.600 & 0.711 & 0.651 & 0.018 \\
\hline
\specialcell{3cm}{По всій \\ вибірці} & 62 & 0.547 & 0.595 & 0.570 & 0.038 \\
\hline
\specialcell{3cm}{По всій \\ вибірці} & 63 & 0.627 & 0.812 & 0.708 & 0.041 \\
\hline
\specialcell{3cm}{Тестовий \\ контекст} & 1 & 1.000 & 0.923 & 0.960 & 0.257 \\
\hline
\specialcell{3cm}{Тестовий \\ контекст} & 2 & 0.839 & 0.959 & 0.895 & 0.485 \\
\hline
\specialcell{3cm}{Тестовий \\ контекст} & 3 & 0.905 & 0.731 & 0.809 & 0.257 \\
\end{mytable*}

\begin{mytable}{ | c | c | c | c | c | c | }%
	{Результати моделювання конкретних реакцій другого набору даних використовуючи ІРС з послідовностями розміром 1--3}%
	{\label{tbl:detailed_data2_irs13}}%
	{№ Контексту & № Реакції & Точність & Прецизійність & Повнота & F-міра}
	
	\specialcell{3cm}{Тестовий \\ контекст} & 1 & 0.151 & 0.052 & 0.077 & 0.272 \\
\hline
\specialcell{3cm}{Тестовий \\ контекст} & 2 & 0.280 & 0.451 & 0.345 & 0.376 \\
\hline
\specialcell{3cm}{Тестовий \\ контекст} & 3 & 0.248 & 0.212 & 0.229 & 0.351 \\
\end{mytable}

\begin{mytable}{ | c | c | c | c | c | c | }%
	{Результати моделювання конкретних реакцій другого набору даних використовуючи ІРС з послідовностями розміром 2--4}%
	{\label{tbl:detailed_data2_irs24}}%
	{№ Контексту & № Реакції & Точність & Прецизійність & Повнота & F-міра}
	
	\specialcell{3cm}{Тестовий \\ контекст} & 1 & 0.752 & 0.385 & 0.509 & 0.272 \\
\hline
\specialcell{3cm}{Тестовий \\ контекст} & 2 & 0.545 & 0.727 & 0.623 & 0.376 \\
\hline
\specialcell{3cm}{Тестовий \\ контекст} & 3 & 0.611 & 0.625 & 0.618 & 0.351 \\
\end{mytable}

\begin{mytable}{ | c | c | c | c | c | c | }%
	{Результати моделювання конкретних реакцій другого набору даних використовуючи ЗНМ з послідовностями розміром 2--4}%
	{\label{tbl:detailed_data2_cnn}}%
	{№ Контексту & № Реакції & Точність & Прецизійність & Повнота & F-міра}
	
	\specialcell{3cm}{Тестовий \\ контекст} & 1 & 0.971 & 0.937 & 0.954 & 0.272 \\
\hline
\specialcell{3cm}{Тестовий \\ контекст} & 2 & 0.939 & 0.925 & 0.932 & 0.376 \\
\hline
\specialcell{3cm}{Тестовий \\ контекст} & 3 & 0.906 & 0.945 & 0.925 & 0.351 \\
\end{mytable}

\FloatBlock

\begin{mytable*}{ | c | c | c | c | c | c | }%
	{Результати моделювання конкретних реакцій третього набору даних використовуючи ІРС з послідовностями розміром 1--3}%
	{\label{tbl:detailed_data3_irs13}}%
	{№ Контексту & № Реакції & Точність & Прецизійність & Повнота & F-міра}
	
	1 & 5 & 0.864 & 0.760 & 0.809 & 0.500 \\
\hline
1 & 6 & 0.786 & 0.880 & 0.830 & 0.500 \\
\hline
3 & 1 & 0.516 & 0.960 & 0.671 & 0.143 \\
\hline
3 & 32 & 0.558 & 0.580 & 0.569 & 0.143 \\
\hline
3 & 33 & 0.854 & 0.700 & 0.769 & 0.143 \\
\hline
3 & 35 & 0.649 & 0.740 & 0.692 & 0.143 \\
\hline
3 & 55 & 0.889 & 0.640 & 0.744 & 0.143 \\
\hline
3 & 56 & 0.556 & 0.100 & 0.169 & 0.143 \\
\hline
3 & 57 & 0.597 & 0.740 & 0.661 & 0.143 \\
\hline
4 & 8 & 0.517 & 0.900 & 0.657 & 0.250 \\
\hline
4 & 31 & 1.000 & 0.440 & 0.611 & 0.250 \\
\hline
4 & 58 & 0.764 & 0.840 & 0.800 & 0.250 \\
\hline
4 & 59 & 0.917 & 0.660 & 0.767 & 0.250 \\
\hline
5 & 8 & 0.441 & 0.820 & 0.573 & 0.167 \\
\hline
5 & 31 & 1.000 & 0.400 & 0.571 & 0.167 \\
\hline
5 & 60 & 0.776 & 0.760 & 0.768 & 0.167 \\
\hline
5 & 61 & 0.610 & 0.500 & 0.549 & 0.167 \\
\hline
5 & 62 & 0.906 & 0.580 & 0.707 & 0.167 \\
\hline
5 & 63 & 0.554 & 0.720 & 0.626 & 0.167 \\
\hline
6 & 8 & 0.473 & 0.860 & 0.610 & 0.250 \\
\hline
6 & 31 & 0.905 & 0.380 & 0.535 & 0.250 \\
\hline
6 & 64 & 0.700 & 0.700 & 0.700 & 0.250 \\
\hline
6 & 65 & 0.737 & 0.560 & 0.636 & 0.250 \\
\hline
7 & 2 & 0.424 & 0.560 & 0.483 & 0.100 \\
\hline
7 & 32 & 0.400 & 0.440 & 0.419 & 0.100 \\
\hline
7 & 33 & 0.840 & 0.420 & 0.560 & 0.100 \\
\hline
7 & 34 & 0.425 & 0.340 & 0.378 & 0.100 \\
\hline
7 & 35 & 0.612 & 0.600 & 0.606 & 0.100 \\
\hline
7 & 36 & 0.622 & 0.560 & 0.589 & 0.100 \\
\hline
7 & 37 & 0.825 & 0.660 & 0.733 & 0.100 \\
\hline
7 & 38 & 0.625 & 0.400 & 0.488 & 0.100 \\
\hline
7 & 39 & 0.478 & 0.640 & 0.547 & 0.100 \\
\hline
7 & 40 & 0.383 & 0.620 & 0.473 & 0.100 \\
\hline
8 & 8 & 0.376 & 0.940 & 0.537 & 0.167 \\
\hline
8 & 31 & 1.000 & 0.340 & 0.507 & 0.167 \\
\hline
8 & 51 & 1.000 & 0.720 & 0.837 & 0.167 \\
\hline
8 & 52 & 0.712 & 0.740 & 0.725 & 0.167 \\
\hline
8 & 53 & 0.957 & 0.900 & 0.928 & 0.167 \\
\hline
8 & 54 & 0.783 & 0.360 & 0.493 & 0.167 \\
\hline
9 & 8 & 0.400 & 0.920 & 0.558 & 0.200 \\
\hline
9 & 31 & 0.909 & 0.400 & 0.556 & 0.200 \\
\hline
9 & 48 & 0.960 & 0.480 & 0.640 & 0.200 \\
\hline
9 & 49 & 0.744 & 0.580 & 0.652 & 0.200 \\
\hline
9 & 50 & 0.939 & 0.920 & 0.929 & 0.200 \\
\hline
10 & 8 & 0.440 & 0.880 & 0.587 & 0.200 \\
\hline
10 & 31 & 1.000 & 0.380 & 0.551 & 0.200 \\
\hline
10 & 42 & 0.809 & 0.760 & 0.784 & 0.200 \\
\hline
10 & 43 & 0.837 & 0.820 & 0.828 & 0.200 \\
\hline
10 & 44 & 0.914 & 0.640 & 0.753 & 0.200 \\
\hline
11 & 8 & 0.430 & 0.800 & 0.559 & 0.200 \\
\hline
11 & 31 & 1.000 & 0.280 & 0.438 & 0.200 \\
\hline
11 & 45 & 0.722 & 0.520 & 0.605 & 0.200 \\
\hline
11 & 46 & 0.569 & 0.660 & 0.611 & 0.200 \\
\hline
11 & 47 & 0.796 & 0.780 & 0.788 & 0.200 \\
\hline
12 & 3 & 0.347 & 0.700 & 0.464 & 0.111 \\
\hline
12 & 7 & 0.485 & 0.320 & 0.386 & 0.111 \\
\hline
12 & 8 & 0.442 & 0.760 & 0.559 & 0.111 \\
\hline
12 & 28 & 0.500 & 0.480 & 0.490 & 0.111 \\
\hline
12 & 32 & 0.444 & 0.400 & 0.421 & 0.111 \\
\hline
12 & 33 & 0.783 & 0.360 & 0.493 & 0.111 \\
\hline
12 & 34 & 0.500 & 0.320 & 0.390 & 0.111 \\
\hline
12 & 35 & 0.676 & 0.500 & 0.575 & 0.111 \\
\hline
12 & 36 & 0.711 & 0.640 & 0.674 & 0.111 \\
\hline
13 & 4 & 0.680 & 0.680 & 0.680 & 0.333 \\
\hline
13 & 5 & 0.667 & 0.600 & 0.632 & 0.333 \\
\hline
13 & 6 & 0.745 & 0.820 & 0.781 & 0.333 \\
\hline
14 & 8 & 0.263 & 0.600 & 0.366 & 0.167 \\
\hline
14 & 9 & 0.462 & 0.240 & 0.316 & 0.167 \\
\hline
14 & 10 & 0.789 & 0.600 & 0.682 & 0.167 \\
\hline
14 & 11 & 0.900 & 0.540 & 0.675 & 0.167 \\
\hline
14 & 12 & 0.410 & 0.640 & 0.500 & 0.167 \\
\hline
14 & 31 & 1.000 & 0.280 & 0.438 & 0.167 \\
\hline
15 & 8 & 0.470 & 0.780 & 0.586 & 0.167 \\
\hline
15 & 13 & 0.589 & 0.660 & 0.623 & 0.167 \\
\hline
15 & 14 & 0.653 & 0.640 & 0.646 & 0.167 \\
\hline
15 & 15 & 0.721 & 0.620 & 0.667 & 0.167 \\
\hline
15 & 16 & 0.554 & 0.620 & 0.585 & 0.167 \\
\hline
15 & 31 & 1.000 & 0.260 & 0.413 & 0.167 \\
\hline
16 & 8 & 0.458 & 0.880 & 0.603 & 0.200 \\
\hline
16 & 17 & 0.769 & 0.800 & 0.784 & 0.200 \\
\hline
16 & 18 & 0.824 & 0.560 & 0.667 & 0.200 \\
\hline
16 & 19 & 0.667 & 0.680 & 0.673 & 0.200 \\
\hline
16 & 31 & 1.000 & 0.340 & 0.507 & 0.200 \\
\hline
17 & 8 & 0.398 & 0.780 & 0.527 & 0.143 \\
\hline
17 & 20 & 0.950 & 0.760 & 0.844 & 0.143 \\
\hline
17 & 21 & 0.311 & 0.280 & 0.295 & 0.143 \\
\hline
17 & 22 & 0.238 & 0.200 & 0.217 & 0.143 \\
\hline
17 & 23 & 0.306 & 0.300 & 0.303 & 0.143 \\
\hline
17 & 24 & 0.322 & 0.380 & 0.349 & 0.143 \\
\hline
17 & 31 & 0.941 & 0.320 & 0.478 & 0.143 \\
\hline
18 & 8 & 0.420 & 0.840 & 0.560 & 0.200 \\
\hline
18 & 25 & 0.720 & 0.360 & 0.480 & 0.200 \\
\hline
18 & 26 & 0.868 & 0.660 & 0.750 & 0.200 \\
\hline
18 & 27 & 0.652 & 0.860 & 0.741 & 0.200 \\
\hline
18 & 31 & 0.857 & 0.360 & 0.507 & 0.200 \\
\hline
19 & 8 & 0.433 & 0.840 & 0.571 & 0.250 \\
\hline
19 & 29 & 0.875 & 0.560 & 0.683 & 0.250 \\
\hline
19 & 30 & 0.596 & 0.620 & 0.608 & 0.250 \\
\hline
19 & 31 & 0.947 & 0.360 & 0.522 & 0.250 \\
\hline
\specialcell{3cm}{По всій \\ вибірці} & 1 & 0.382 & 0.260 & 0.310 & 0.016 \\
\hline
\specialcell{3cm}{По всій \\ вибірці} & 2 & 0.152 & 0.280 & 0.197 & 0.016 \\
\hline
\specialcell{3cm}{По всій \\ вибірці} & 3 & 0.132 & 0.240 & 0.170 & 0.016 \\
\hline
\specialcell{3cm}{По всій \\ вибірці} & 4 & 0.081 & 0.220 & 0.119 & 0.016 \\
\hline
\specialcell{3cm}{По всій \\ вибірці} & 5 & 0.189 & 0.280 & 0.226 & 0.016 \\
\hline
\specialcell{3cm}{По всій \\ вибірці} & 6 & 0.810 & 0.340 & 0.479 & 0.016 \\
\hline
\specialcell{3cm}{По всій \\ вибірці} & 7 & 0.300 & 0.120 & 0.171 & 0.016 \\
\hline
\specialcell{3cm}{По всій \\ вибірці} & 8 & 0.112 & 0.460 & 0.180 & 0.016 \\
\hline
\specialcell{3cm}{По всій \\ вибірці} & 9 & 0.111 & 0.140 & 0.124 & 0.016 \\
\hline
\specialcell{3cm}{По всій \\ вибірці} & 10 & 0.202 & 0.360 & 0.259 & 0.016 \\
\hline
\specialcell{3cm}{По всій \\ вибірці} & 11 & 0.235 & 0.080 & 0.119 & 0.016 \\
\hline
\specialcell{3cm}{По всій \\ вибірці} & 12 & 0.150 & 0.460 & 0.227 & 0.016 \\
\hline
\specialcell{3cm}{По всій \\ вибірці} & 13 & 0.350 & 0.420 & 0.382 & 0.016 \\
\hline
\specialcell{3cm}{По всій \\ вибірці} & 14 & 0.415 & 0.440 & 0.427 & 0.016 \\
\hline
\specialcell{3cm}{По всій \\ вибірці} & 15 & 0.286 & 0.360 & 0.319 & 0.016 \\
\hline
\specialcell{3cm}{По всій \\ вибірці} & 16 & 0.194 & 0.360 & 0.252 & 0.016 \\
\hline
\specialcell{3cm}{По всій \\ вибірці} & 17 & 0.513 & 0.400 & 0.449 & 0.016 \\
\hline
\specialcell{3cm}{По всій \\ вибірці} & 18 & 0.667 & 0.280 & 0.394 & 0.016 \\
\hline
\specialcell{3cm}{По всій \\ вибірці} & 19 & 0.255 & 0.240 & 0.247 & 0.016 \\
\hline
\specialcell{3cm}{По всій \\ вибірці} & 20 & 0.250 & 0.360 & 0.295 & 0.016 \\
\hline
\specialcell{3cm}{По всій \\ вибірці} & 21 & 0.297 & 0.220 & 0.253 & 0.016 \\
\hline
\specialcell{3cm}{По всій \\ вибірці} & 22 & 0.205 & 0.160 & 0.180 & 0.016 \\
\hline
\specialcell{3cm}{По всій \\ вибірці} & 23 & 0.161 & 0.180 & 0.170 & 0.016 \\
\hline
\specialcell{3cm}{По всій \\ вибірці} & 24 & 0.217 & 0.200 & 0.208 & 0.016 \\
\hline
\specialcell{3cm}{По всій \\ вибірці} & 25 & 0.375 & 0.180 & 0.243 & 0.016 \\
\hline
\specialcell{3cm}{По всій \\ вибірці} & 26 & 0.385 & 0.300 & 0.337 & 0.016 \\
\hline
\specialcell{3cm}{По всій \\ вибірці} & 27 & 0.344 & 0.420 & 0.378 & 0.016 \\
\hline
\specialcell{3cm}{По всій \\ вибірці} & 28 & 0.218 & 0.240 & 0.229 & 0.016 \\
\hline
\specialcell{3cm}{По всій \\ вибірці} & 29 & 0.261 & 0.240 & 0.250 & 0.016 \\
\hline
\specialcell{3cm}{По всій \\ вибірці} & 30 & 0.245 & 0.240 & 0.242 & 0.016 \\
\hline
\specialcell{3cm}{По всій \\ вибірці} & 31 & 1.000 & 0.200 & 0.333 & 0.016 \\
\hline
\specialcell{3cm}{По всій \\ вибірці} & 32 & 0.227 & 0.200 & 0.213 & 0.016 \\
\hline
\specialcell{3cm}{По всій \\ вибірці} & 33 & 0.524 & 0.220 & 0.310 & 0.016 \\
\hline
\specialcell{3cm}{По всій \\ вибірці} & 34 & 0.314 & 0.220 & 0.259 & 0.016 \\
\hline
\specialcell{3cm}{По всій \\ вибірці} & 35 & 0.722 & 0.260 & 0.382 & 0.016 \\
\hline
\specialcell{3cm}{По всій \\ вибірці} & 36 & 0.246 & 0.300 & 0.270 & 0.016 \\
\hline
\specialcell{3cm}{По всій \\ вибірці} & 37 & 0.800 & 0.480 & 0.600 & 0.016 \\
\hline
\specialcell{3cm}{По всій \\ вибірці} & 38 & 0.300 & 0.120 & 0.171 & 0.016 \\
\hline
\specialcell{3cm}{По всій \\ вибірці} & 39 & 0.373 & 0.440 & 0.404 & 0.016 \\
\hline
\specialcell{3cm}{По всій \\ вибірці} & 40 & 0.175 & 0.340 & 0.231 & 0.016 \\
\hline
\specialcell{3cm}{По всій \\ вибірці} & 42 & 0.484 & 0.300 & 0.370 & 0.016 \\
\hline
\specialcell{3cm}{По всій \\ вибірці} & 43 & 0.714 & 0.500 & 0.588 & 0.016 \\
\hline
\specialcell{3cm}{По всій \\ вибірці} & 44 & 0.656 & 0.420 & 0.512 & 0.016 \\
\hline
\specialcell{3cm}{По всій \\ вибірці} & 45 & 0.500 & 0.360 & 0.419 & 0.016 \\
\hline
\specialcell{3cm}{По всій \\ вибірці} & 46 & 0.121 & 0.340 & 0.179 & 0.016 \\
\hline
\specialcell{3cm}{По всій \\ вибірці} & 47 & 0.674 & 0.620 & 0.646 & 0.016 \\
\hline
\specialcell{3cm}{По всій \\ вибірці} & 48 & 0.261 & 0.120 & 0.164 & 0.016 \\
\hline
\specialcell{3cm}{По всій \\ вибірці} & 49 & 0.429 & 0.120 & 0.188 & 0.016 \\
\hline
\specialcell{3cm}{По всій \\ вибірці} & 50 & 0.879 & 0.580 & 0.699 & 0.016 \\
\hline
\specialcell{3cm}{По всій \\ вибірці} & 51 & 1.000 & 0.400 & 0.571 & 0.016 \\
\hline
\specialcell{3cm}{По всій \\ вибірці} & 52 & 0.667 & 0.360 & 0.468 & 0.016 \\
\hline
\specialcell{3cm}{По всій \\ вибірці} & 53 & 0.578 & 0.520 & 0.547 & 0.016 \\
\hline
\specialcell{3cm}{По всій \\ вибірці} & 54 & 0.467 & 0.140 & 0.215 & 0.016 \\
\hline
\specialcell{3cm}{По всій \\ вибірці} & 55 & 0.684 & 0.260 & 0.377 & 0.016 \\
\hline
\specialcell{3cm}{По всій \\ вибірці} & 56 & 1.000 & 0.020 & 0.039 & 0.016 \\
\hline
\specialcell{3cm}{По всій \\ вибірці} & 57 & 0.333 & 0.100 & 0.154 & 0.016 \\
\hline
\specialcell{3cm}{По всій \\ вибірці} & 58 & 0.521 & 0.500 & 0.510 & 0.016 \\
\hline
\specialcell{3cm}{По всій \\ вибірці} & 59 & 0.457 & 0.320 & 0.376 & 0.016 \\
\hline
\specialcell{3cm}{По всій \\ вибірці} & 60 & 0.561 & 0.460 & 0.505 & 0.016 \\
\hline
\specialcell{3cm}{По всій \\ вибірці} & 61 & 0.327 & 0.340 & 0.333 & 0.016 \\
\hline
\specialcell{3cm}{По всій \\ вибірці} & 62 & 0.270 & 0.200 & 0.230 & 0.016 \\
\hline
\specialcell{3cm}{По всій \\ вибірці} & 63 & 0.308 & 0.480 & 0.375 & 0.016 \\
\hline
\specialcell{3cm}{По всій \\ вибірці} & 64 & 0.279 & 0.240 & 0.258 & 0.016 \\
\hline
\specialcell{3cm}{По всій \\ вибірці} & 65 & 0.318 & 0.280 & 0.298 & 0.016 \\
\hline
\specialcell{3cm}{Тестовий \\ контекст} & 1 & 0.822 & 0.740 & 0.779 & 0.333 \\
\hline
\specialcell{3cm}{Тестовий \\ контекст} & 2 & 0.712 & 0.740 & 0.725 & 0.333 \\
\hline
\specialcell{3cm}{Тестовий \\ контекст} & 3 & 0.679 & 0.720 & 0.699 & 0.333 \\
\end{mytable*}

\begin{mytable*}{ | c | c | c | c | c | c | }%
	{Результати моделювання конкретних реакцій третього набору даних використовуючи ІРС з послідовностями розміром 2--4}%
	{\label{tbl:detailed_data3_irs24}}%
	{№ Контексту & № Реакції & Точність & Прецизійність & Повнота & F-міра}
	
	1 & 5 & 0.918 & 0.900 & 0.909 & 0.500 \\
\hline
1 & 6 & 0.930 & 0.800 & 0.860 & 0.500 \\
\hline
3 & 1 & 0.731 & 0.980 & 0.838 & 0.143 \\
\hline
3 & 32 & 0.804 & 0.900 & 0.849 & 0.143 \\
\hline
3 & 33 & 1.000 & 0.800 & 0.889 & 0.143 \\
\hline
3 & 35 & 0.941 & 0.960 & 0.950 & 0.143 \\
\hline
3 & 55 & 0.979 & 0.940 & 0.959 & 0.143 \\
\hline
3 & 56 & 0.931 & 0.540 & 0.684 & 0.143 \\
\hline
3 & 57 & 0.797 & 0.940 & 0.862 & 0.143 \\
\hline
4 & 8 & 0.687 & 0.920 & 0.786 & 0.250 \\
\hline
4 & 31 & 1.000 & 0.600 & 0.750 & 0.250 \\
\hline
4 & 58 & 0.926 & 1.000 & 0.962 & 0.250 \\
\hline
4 & 59 & 0.980 & 0.960 & 0.970 & 0.250 \\
\hline
5 & 8 & 0.600 & 0.900 & 0.720 & 0.167 \\
\hline
5 & 31 & 1.000 & 0.560 & 0.718 & 0.167 \\
\hline
5 & 60 & 0.880 & 0.880 & 0.880 & 0.167 \\
\hline
5 & 61 & 0.851 & 0.800 & 0.825 & 0.167 \\
\hline
5 & 62 & 0.956 & 0.860 & 0.905 & 0.167 \\
\hline
5 & 63 & 0.891 & 0.980 & 0.933 & 0.167 \\
\hline
6 & 8 & 0.662 & 0.980 & 0.790 & 0.250 \\
\hline
6 & 31 & 1.000 & 0.520 & 0.684 & 0.250 \\
\hline
6 & 64 & 0.904 & 0.940 & 0.922 & 0.250 \\
\hline
6 & 65 & 0.938 & 0.900 & 0.918 & 0.250 \\
\hline
7 & 2 & 0.719 & 0.820 & 0.766 & 0.100 \\
\hline
7 & 32 & 0.750 & 0.720 & 0.735 & 0.100 \\
\hline
7 & 33 & 0.914 & 0.640 & 0.753 & 0.100 \\
\hline
7 & 34 & 0.816 & 0.800 & 0.808 & 0.100 \\
\hline
7 & 35 & 0.820 & 0.820 & 0.820 & 0.100 \\
\hline
7 & 36 & 0.778 & 0.700 & 0.737 & 0.100 \\
\hline
7 & 37 & 0.955 & 0.840 & 0.894 & 0.100 \\
\hline
7 & 38 & 0.912 & 0.620 & 0.738 & 0.100 \\
\hline
7 & 39 & 0.676 & 0.920 & 0.780 & 0.100 \\
\hline
7 & 40 & 0.614 & 0.860 & 0.717 & 0.100 \\
\hline
8 & 8 & 0.698 & 0.880 & 0.779 & 0.167 \\
\hline
8 & 31 & 1.000 & 0.520 & 0.684 & 0.167 \\
\hline
8 & 51 & 0.962 & 1.000 & 0.980 & 0.167 \\
\hline
8 & 52 & 0.738 & 0.960 & 0.835 & 0.167 \\
\hline
8 & 53 & 0.961 & 0.980 & 0.970 & 0.167 \\
\hline
8 & 54 & 0.930 & 0.800 & 0.860 & 0.167 \\
\hline
9 & 8 & 0.676 & 0.960 & 0.793 & 0.200 \\
\hline
9 & 31 & 1.000 & 0.600 & 0.750 & 0.200 \\
\hline
9 & 48 & 1.000 & 0.920 & 0.958 & 0.200 \\
\hline
9 & 49 & 0.941 & 0.960 & 0.950 & 0.200 \\
\hline
9 & 50 & 0.961 & 0.980 & 0.970 & 0.200 \\
\hline
10 & 8 & 0.634 & 0.900 & 0.744 & 0.200 \\
\hline
10 & 31 & 1.000 & 0.600 & 0.750 & 0.200 \\
\hline
10 & 42 & 0.922 & 0.940 & 0.931 & 0.200 \\
\hline
10 & 43 & 0.961 & 0.980 & 0.970 & 0.200 \\
\hline
10 & 44 & 0.957 & 0.900 & 0.928 & 0.200 \\
\hline
11 & 8 & 0.708 & 0.920 & 0.800 & 0.200 \\
\hline
11 & 31 & 1.000 & 0.460 & 0.630 & 0.200 \\
\hline
11 & 45 & 0.776 & 0.900 & 0.833 & 0.200 \\
\hline
11 & 46 & 0.818 & 0.900 & 0.857 & 0.200 \\
\hline
11 & 47 & 0.959 & 0.940 & 0.949 & 0.200 \\
\hline
12 & 3 & 0.500 & 0.900 & 0.643 & 0.111 \\
\hline
12 & 7 & 0.750 & 0.480 & 0.585 & 0.111 \\
\hline
12 & 8 & 0.827 & 0.860 & 0.843 & 0.111 \\
\hline
12 & 28 & 0.894 & 0.840 & 0.866 & 0.111 \\
\hline
12 & 32 & 0.759 & 0.820 & 0.788 & 0.111 \\
\hline
12 & 33 & 0.914 & 0.640 & 0.753 & 0.111 \\
\hline
12 & 34 & 0.830 & 0.780 & 0.804 & 0.111 \\
\hline
12 & 35 & 0.771 & 0.740 & 0.755 & 0.111 \\
\hline
12 & 36 & 0.778 & 0.700 & 0.737 & 0.111 \\
\hline
13 & 4 & 0.721 & 0.980 & 0.831 & 0.333 \\
\hline
13 & 5 & 0.973 & 0.720 & 0.828 & 0.333 \\
\hline
13 & 6 & 0.925 & 0.740 & 0.822 & 0.333 \\
\hline
14 & 8 & 0.488 & 0.820 & 0.612 & 0.167 \\
\hline
14 & 9 & 0.674 & 0.620 & 0.646 & 0.167 \\
\hline
14 & 10 & 0.870 & 0.800 & 0.833 & 0.167 \\
\hline
14 & 11 & 0.952 & 0.800 & 0.870 & 0.167 \\
\hline
14 & 12 & 0.741 & 0.860 & 0.796 & 0.167 \\
\hline
14 & 31 & 1.000 & 0.480 & 0.649 & 0.167 \\
\hline
15 & 8 & 0.553 & 0.840 & 0.667 & 0.167 \\
\hline
15 & 13 & 0.868 & 0.920 & 0.893 & 0.167 \\
\hline
15 & 14 & 0.854 & 0.820 & 0.837 & 0.167 \\
\hline
15 & 15 & 0.932 & 0.820 & 0.872 & 0.167 \\
\hline
15 & 16 & 0.737 & 0.840 & 0.785 & 0.167 \\
\hline
15 & 31 & 1.000 & 0.440 & 0.611 & 0.167 \\
\hline
16 & 8 & 0.677 & 0.880 & 0.765 & 0.200 \\
\hline
16 & 17 & 0.868 & 0.920 & 0.893 & 0.200 \\
\hline
16 & 18 & 0.941 & 0.960 & 0.950 & 0.200 \\
\hline
16 & 19 & 0.839 & 0.940 & 0.887 & 0.200 \\
\hline
16 & 31 & 1.000 & 0.500 & 0.667 & 0.200 \\
\hline
17 & 8 & 0.551 & 0.860 & 0.672 & 0.143 \\
\hline
17 & 20 & 0.935 & 0.860 & 0.896 & 0.143 \\
\hline
17 & 21 & 0.577 & 0.600 & 0.588 & 0.143 \\
\hline
17 & 22 & 0.562 & 0.720 & 0.632 & 0.143 \\
\hline
17 & 23 & 0.522 & 0.480 & 0.500 & 0.143 \\
\hline
17 & 24 & 0.696 & 0.640 & 0.667 & 0.143 \\
\hline
17 & 31 & 1.000 & 0.360 & 0.529 & 0.143 \\
\hline
18 & 8 & 0.512 & 0.880 & 0.647 & 0.200 \\
\hline
18 & 25 & 0.750 & 0.540 & 0.628 & 0.200 \\
\hline
18 & 26 & 0.952 & 0.800 & 0.870 & 0.200 \\
\hline
18 & 27 & 0.803 & 0.980 & 0.883 & 0.200 \\
\hline
18 & 31 & 1.000 & 0.500 & 0.667 & 0.200 \\
\hline
19 & 8 & 0.667 & 0.840 & 0.743 & 0.250 \\
\hline
19 & 29 & 0.902 & 0.920 & 0.911 & 0.250 \\
\hline
19 & 30 & 0.842 & 0.960 & 0.897 & 0.250 \\
\hline
19 & 31 & 1.000 & 0.580 & 0.734 & 0.250 \\
\hline
\specialcell{3cm}{По всій \\ вибірці} & 1 & 0.677 & 0.840 & 0.750 & 0.016 \\
\hline
\specialcell{3cm}{По всій \\ вибірці} & 2 & 0.423 & 0.600 & 0.496 & 0.016 \\
\hline
\specialcell{3cm}{По всій \\ вибірці} & 3 & 0.326 & 0.620 & 0.428 & 0.016 \\
\hline
\specialcell{3cm}{По всій \\ вибірці} & 4 & 0.429 & 0.540 & 0.478 & 0.016 \\
\hline
\specialcell{3cm}{По всій \\ вибірці} & 5 & 0.789 & 0.300 & 0.435 & 0.016 \\
\hline
\specialcell{3cm}{По всій \\ вибірці} & 6 & 1.000 & 0.320 & 0.485 & 0.016 \\
\hline
\specialcell{3cm}{По всій \\ вибірці} & 7 & 0.762 & 0.320 & 0.451 & 0.016 \\
\hline
\specialcell{3cm}{По всій \\ вибірці} & 8 & 0.307 & 0.620 & 0.411 & 0.016 \\
\hline
\specialcell{3cm}{По всій \\ вибірці} & 9 & 0.385 & 0.500 & 0.435 & 0.016 \\
\hline
\specialcell{3cm}{По всій \\ вибірці} & 10 & 0.403 & 0.540 & 0.462 & 0.016 \\
\hline
\specialcell{3cm}{По всій \\ вибірці} & 11 & 0.645 & 0.400 & 0.494 & 0.016 \\
\hline
\specialcell{3cm}{По всій \\ вибірці} & 12 & 0.427 & 0.700 & 0.530 & 0.016 \\
\hline
\specialcell{3cm}{По всій \\ вибірці} & 13 & 0.791 & 0.680 & 0.731 & 0.016 \\
\hline
\specialcell{3cm}{По всій \\ вибірці} & 14 & 0.865 & 0.640 & 0.736 & 0.016 \\
\hline
\specialcell{3cm}{По всій \\ вибірці} & 15 & 0.560 & 0.560 & 0.560 & 0.016 \\
\hline
\specialcell{3cm}{По всій \\ вибірці} & 16 & 0.461 & 0.700 & 0.556 & 0.016 \\
\hline
\specialcell{3cm}{По всій \\ вибірці} & 17 & 0.895 & 0.680 & 0.773 & 0.016 \\
\hline
\specialcell{3cm}{По всій \\ вибірці} & 18 & 0.714 & 0.900 & 0.796 & 0.016 \\
\hline
\specialcell{3cm}{По всій \\ вибірці} & 19 & 0.679 & 0.760 & 0.717 & 0.016 \\
\hline
\specialcell{3cm}{По всій \\ вибірці} & 20 & 0.574 & 0.700 & 0.631 & 0.016 \\
\hline
\specialcell{3cm}{По всій \\ вибірці} & 21 & 0.596 & 0.560 & 0.577 & 0.016 \\
\hline
\specialcell{3cm}{По всій \\ вибірці} & 22 & 0.569 & 0.660 & 0.611 & 0.016 \\
\hline
\specialcell{3cm}{По всій \\ вибірці} & 23 & 0.395 & 0.340 & 0.366 & 0.016 \\
\hline
\specialcell{3cm}{По всій \\ вибірці} & 24 & 0.650 & 0.520 & 0.578 & 0.016 \\
\hline
\specialcell{3cm}{По всій \\ вибірці} & 25 & 0.857 & 0.240 & 0.375 & 0.016 \\
\hline
\specialcell{3cm}{По всій \\ вибірці} & 26 & 0.833 & 0.500 & 0.625 & 0.016 \\
\hline
\specialcell{3cm}{По всій \\ вибірці} & 27 & 0.623 & 0.660 & 0.641 & 0.016 \\
\hline
\specialcell{3cm}{По всій \\ вибірці} & 28 & 0.509 & 0.580 & 0.542 & 0.016 \\
\hline
\specialcell{3cm}{По всій \\ вибірці} & 29 & 0.538 & 0.560 & 0.549 & 0.016 \\
\hline
\specialcell{3cm}{По всій \\ вибірці} & 30 & 0.638 & 0.880 & 0.739 & 0.016 \\
\hline
\specialcell{3cm}{По всій \\ вибірці} & 31 & 1.000 & 0.300 & 0.462 & 0.016 \\
\hline
\specialcell{3cm}{По всій \\ вибірці} & 32 & 0.674 & 0.580 & 0.624 & 0.016 \\
\hline
\specialcell{3cm}{По всій \\ вибірці} & 33 & 0.857 & 0.600 & 0.706 & 0.016 \\
\hline
\specialcell{3cm}{По всій \\ вибірці} & 34 & 0.767 & 0.660 & 0.710 & 0.016 \\
\hline
\specialcell{3cm}{По всій \\ вибірці} & 35 & 0.778 & 0.560 & 0.651 & 0.016 \\
\hline
\specialcell{3cm}{По всій \\ вибірці} & 36 & 0.528 & 0.560 & 0.544 & 0.016 \\
\hline
\specialcell{3cm}{По всій \\ вибірці} & 37 & 0.868 & 0.660 & 0.750 & 0.016 \\
\hline
\specialcell{3cm}{По всій \\ вибірці} & 38 & 0.750 & 0.480 & 0.585 & 0.016 \\
\hline
\specialcell{3cm}{По всій \\ вибірці} & 39 & 0.554 & 0.820 & 0.661 & 0.016 \\
\hline
\specialcell{3cm}{По всій \\ вибірці} & 40 & 0.371 & 0.660 & 0.475 & 0.016 \\
\hline
\specialcell{3cm}{По всій \\ вибірці} & 42 & 0.727 & 0.480 & 0.578 & 0.016 \\
\hline
\specialcell{3cm}{По всій \\ вибірці} & 43 & 0.836 & 0.920 & 0.876 & 0.016 \\
\hline
\specialcell{3cm}{По всій \\ вибірці} & 44 & 0.886 & 0.780 & 0.830 & 0.016 \\
\hline
\specialcell{3cm}{По всій \\ вибірці} & 45 & 0.735 & 0.720 & 0.727 & 0.016 \\
\hline
\specialcell{3cm}{По всій \\ вибірці} & 46 & 0.364 & 0.560 & 0.441 & 0.016 \\
\hline
\specialcell{3cm}{По всій \\ вибірці} & 47 & 0.854 & 0.820 & 0.837 & 0.016 \\
\hline
\specialcell{3cm}{По всій \\ вибірці} & 48 & 0.783 & 0.720 & 0.750 & 0.016 \\
\hline
\specialcell{3cm}{По всій \\ вибірці} & 49 & 0.824 & 0.560 & 0.667 & 0.016 \\
\hline
\specialcell{3cm}{По всій \\ вибірці} & 50 & 0.868 & 0.920 & 0.893 & 0.016 \\
\hline
\specialcell{3cm}{По всій \\ вибірці} & 51 & 0.889 & 0.960 & 0.923 & 0.016 \\
\hline
\specialcell{3cm}{По всій \\ вибірці} & 52 & 0.754 & 0.860 & 0.804 & 0.016 \\
\hline
\specialcell{3cm}{По всій \\ вибірці} & 53 & 0.961 & 0.980 & 0.970 & 0.016 \\
\hline
\specialcell{3cm}{По всій \\ вибірці} & 54 & 0.674 & 0.580 & 0.624 & 0.016 \\
\hline
\specialcell{3cm}{По всій \\ вибірці} & 55 & 0.841 & 0.740 & 0.787 & 0.016 \\
\hline
\specialcell{3cm}{По всій \\ вибірці} & 56 & 0.846 & 0.440 & 0.579 & 0.016 \\
\hline
\specialcell{3cm}{По всій \\ вибірці} & 57 & 0.692 & 0.720 & 0.706 & 0.016 \\
\hline
\specialcell{3cm}{По всій \\ вибірці} & 58 & 0.860 & 0.860 & 0.860 & 0.016 \\
\hline
\specialcell{3cm}{По всій \\ вибірці} & 59 & 0.905 & 0.760 & 0.826 & 0.016 \\
\hline
\specialcell{3cm}{По всій \\ вибірці} & 60 & 0.860 & 0.740 & 0.796 & 0.016 \\
\hline
\specialcell{3cm}{По всій \\ вибірці} & 61 & 0.654 & 0.680 & 0.667 & 0.016 \\
\hline
\specialcell{3cm}{По всій \\ вибірці} & 62 & 0.711 & 0.540 & 0.614 & 0.016 \\
\hline
\specialcell{3cm}{По всій \\ вибірці} & 63 & 0.541 & 0.920 & 0.681 & 0.016 \\
\hline
\specialcell{3cm}{По всій \\ вибірці} & 64 & 0.514 & 0.380 & 0.437 & 0.016 \\
\hline
\specialcell{3cm}{По всій \\ вибірці} & 65 & 0.872 & 0.820 & 0.845 & 0.016 \\
\hline
\specialcell{3cm}{Тестовий \\ контекст} & 1 & 0.979 & 0.940 & 0.959 & 0.333 \\
\hline
\specialcell{3cm}{Тестовий \\ контекст} & 2 & 0.849 & 0.900 & 0.874 & 0.333 \\
\hline
\specialcell{3cm}{Тестовий \\ контекст} & 3 & 0.857 & 0.840 & 0.848 & 0.333 \\
\end{mytable*}

\begin{mytable*}{ | c | c | c | c | c | c | }%
	{Результати моделювання конкретних реакцій третього набору даних використовуючи ЗНМ з послідовностями розміром 2--4}%
	{\label{tbl:detailed_data3_cnn}}%
	{№ Контексту & № Реакції & Точність & Прецизійність & Повнота & F-міра}
	
	1 & 5 & 0.885 & 0.920 & 0.902 & 0.500 \\
\hline
1 & 6 & 0.917 & 0.880 & 0.898 & 0.500 \\
\hline
3 & 1 & 1.000 & 0.980 & 0.990 & 0.143 \\
\hline
3 & 32 & 1.000 & 1.000 & 1.000 & 0.143 \\
\hline
3 & 33 & 1.000 & 1.000 & 1.000 & 0.143 \\
\hline
3 & 35 & 1.000 & 1.000 & 1.000 & 0.143 \\
\hline
3 & 55 & 1.000 & 1.000 & 1.000 & 0.143 \\
\hline
3 & 56 & 1.000 & 1.000 & 1.000 & 0.143 \\
\hline
3 & 57 & 0.980 & 1.000 & 0.990 & 0.143 \\
\hline
4 & 8 & 1.000 & 0.960 & 0.980 & 0.250 \\
\hline
4 & 31 & 1.000 & 1.000 & 1.000 & 0.250 \\
\hline
4 & 58 & 0.980 & 1.000 & 0.990 & 0.250 \\
\hline
4 & 59 & 0.980 & 1.000 & 0.990 & 0.250 \\
\hline
5 & 8 & 0.978 & 0.900 & 0.938 & 0.167 \\
\hline
5 & 31 & 0.980 & 1.000 & 0.990 & 0.167 \\
\hline
5 & 60 & 0.960 & 0.960 & 0.960 & 0.167 \\
\hline
5 & 61 & 0.980 & 0.980 & 0.980 & 0.167 \\
\hline
5 & 62 & 0.960 & 0.960 & 0.960 & 0.167 \\
\hline
5 & 63 & 0.943 & 1.000 & 0.971 & 0.167 \\
\hline
6 & 8 & 1.000 & 0.920 & 0.958 & 0.250 \\
\hline
6 & 31 & 0.980 & 1.000 & 0.990 & 0.250 \\
\hline
6 & 64 & 0.962 & 1.000 & 0.980 & 0.250 \\
\hline
6 & 65 & 0.961 & 0.980 & 0.970 & 0.250 \\
\hline
7 & 2 & 0.815 & 0.880 & 0.846 & 0.100 \\
\hline
7 & 32 & 0.979 & 0.940 & 0.959 & 0.100 \\
\hline
7 & 33 & 0.980 & 1.000 & 0.990 & 0.100 \\
\hline
7 & 34 & 0.980 & 0.960 & 0.970 & 0.100 \\
\hline
7 & 35 & 0.926 & 1.000 & 0.962 & 0.100 \\
\hline
7 & 36 & 0.961 & 0.980 & 0.970 & 0.100 \\
\hline
7 & 37 & 1.000 & 0.980 & 0.990 & 0.100 \\
\hline
7 & 38 & 0.957 & 0.900 & 0.928 & 0.100 \\
\hline
7 & 39 & 0.958 & 0.920 & 0.939 & 0.100 \\
\hline
7 & 40 & 0.939 & 0.920 & 0.929 & 0.100 \\
\hline
8 & 8 & 0.979 & 0.920 & 0.948 & 0.167 \\
\hline
8 & 31 & 0.980 & 1.000 & 0.990 & 0.167 \\
\hline
8 & 51 & 1.000 & 0.980 & 0.990 & 0.167 \\
\hline
8 & 52 & 0.962 & 1.000 & 0.980 & 0.167 \\
\hline
8 & 53 & 0.980 & 1.000 & 0.990 & 0.167 \\
\hline
8 & 54 & 1.000 & 1.000 & 1.000 & 0.167 \\
\hline
9 & 8 & 0.959 & 0.940 & 0.949 & 0.200 \\
\hline
9 & 31 & 0.980 & 1.000 & 0.990 & 0.200 \\
\hline
9 & 48 & 0.959 & 0.940 & 0.949 & 0.200 \\
\hline
9 & 49 & 1.000 & 1.000 & 1.000 & 0.200 \\
\hline
9 & 50 & 0.980 & 1.000 & 0.990 & 0.200 \\
\hline
10 & 8 & 0.978 & 0.900 & 0.938 & 0.200 \\
\hline
10 & 31 & 0.980 & 0.980 & 0.980 & 0.200 \\
\hline
10 & 42 & 0.979 & 0.940 & 0.959 & 0.200 \\
\hline
10 & 43 & 0.926 & 1.000 & 0.962 & 0.200 \\
\hline
10 & 44 & 0.942 & 0.980 & 0.961 & 0.200 \\
\hline
11 & 8 & 1.000 & 0.820 & 0.901 & 0.200 \\
\hline
11 & 31 & 1.000 & 1.000 & 1.000 & 0.200 \\
\hline
11 & 45 & 0.980 & 0.960 & 0.970 & 0.200 \\
\hline
11 & 46 & 0.891 & 0.980 & 0.933 & 0.200 \\
\hline
11 & 47 & 0.909 & 1.000 & 0.952 & 0.200 \\
\hline
12 & 3 & 0.887 & 0.940 & 0.913 & 0.111 \\
\hline
12 & 7 & 0.955 & 0.840 & 0.894 & 0.111 \\
\hline
12 & 8 & 0.957 & 0.880 & 0.917 & 0.111 \\
\hline
12 & 28 & 0.926 & 1.000 & 0.962 & 0.111 \\
\hline
12 & 32 & 0.958 & 0.920 & 0.939 & 0.111 \\
\hline
12 & 33 & 0.980 & 1.000 & 0.990 & 0.111 \\
\hline
12 & 34 & 0.980 & 1.000 & 0.990 & 0.111 \\
\hline
12 & 35 & 0.942 & 0.980 & 0.961 & 0.111 \\
\hline
12 & 36 & 0.980 & 1.000 & 0.990 & 0.111 \\
\hline
13 & 4 & 0.962 & 1.000 & 0.980 & 0.333 \\
\hline
13 & 5 & 0.898 & 0.880 & 0.889 & 0.333 \\
\hline
13 & 6 & 0.918 & 0.900 & 0.909 & 0.333 \\
\hline
14 & 8 & 0.938 & 0.900 & 0.918 & 0.167 \\
\hline
14 & 9 & 0.920 & 0.920 & 0.920 & 0.167 \\
\hline
14 & 10 & 0.980 & 1.000 & 0.990 & 0.167 \\
\hline
14 & 11 & 0.980 & 0.980 & 0.980 & 0.167 \\
\hline
14 & 12 & 0.900 & 0.900 & 0.900 & 0.167 \\
\hline
14 & 31 & 0.980 & 1.000 & 0.990 & 0.167 \\
\hline
15 & 8 & 0.952 & 0.800 & 0.870 & 0.167 \\
\hline
15 & 13 & 0.980 & 1.000 & 0.990 & 0.167 \\
\hline
15 & 14 & 0.922 & 0.940 & 0.931 & 0.167 \\
\hline
15 & 15 & 0.842 & 0.960 & 0.897 & 0.167 \\
\hline
15 & 16 & 0.875 & 0.840 & 0.857 & 0.167 \\
\hline
15 & 31 & 0.980 & 1.000 & 0.990 & 0.167 \\
\hline
16 & 8 & 1.000 & 0.880 & 0.936 & 0.200 \\
\hline
16 & 17 & 0.893 & 1.000 & 0.943 & 0.200 \\
\hline
16 & 18 & 0.962 & 1.000 & 0.980 & 0.200 \\
\hline
16 & 19 & 1.000 & 0.980 & 0.990 & 0.200 \\
\hline
16 & 31 & 1.000 & 0.980 & 0.990 & 0.200 \\
\hline
17 & 8 & 1.000 & 0.880 & 0.936 & 0.143 \\
\hline
17 & 20 & 1.000 & 1.000 & 1.000 & 0.143 \\
\hline
17 & 21 & 0.980 & 0.980 & 0.980 & 0.143 \\
\hline
17 & 22 & 0.885 & 0.920 & 0.902 & 0.143 \\
\hline
17 & 23 & 0.855 & 0.940 & 0.895 & 0.143 \\
\hline
17 & 24 & 0.959 & 0.940 & 0.949 & 0.143 \\
\hline
17 & 31 & 1.000 & 1.000 & 1.000 & 0.143 \\
\hline
18 & 8 & 0.920 & 0.920 & 0.920 & 0.200 \\
\hline
18 & 25 & 0.939 & 0.920 & 0.929 & 0.200 \\
\hline
18 & 26 & 0.922 & 0.940 & 0.931 & 0.200 \\
\hline
18 & 27 & 0.979 & 0.940 & 0.959 & 0.200 \\
\hline
18 & 31 & 0.885 & 0.920 & 0.902 & 0.200 \\
\hline
19 & 8 & 1.000 & 0.920 & 0.958 & 0.250 \\
\hline
19 & 29 & 0.962 & 1.000 & 0.980 & 0.250 \\
\hline
19 & 30 & 0.980 & 1.000 & 0.990 & 0.250 \\
\hline
19 & 31 & 0.980 & 1.000 & 0.990 & 0.250 \\
\hline
\specialcell{3cm}{По всій \\ вибірці} & 1 & 0.776 & 0.760 & 0.768 & 0.016 \\
\hline
\specialcell{3cm}{По всій \\ вибірці} & 2 & 0.827 & 0.860 & 0.843 & 0.016 \\
\hline
\specialcell{3cm}{По всій \\ вибірці} & 3 & 0.872 & 0.820 & 0.845 & 0.016 \\
\hline
\specialcell{3cm}{По всій \\ вибірці} & 4 & 0.786 & 0.660 & 0.717 & 0.016 \\
\hline
\specialcell{3cm}{По всій \\ вибірці} & 5 & 0.860 & 0.740 & 0.796 & 0.016 \\
\hline
\specialcell{3cm}{По всій \\ вибірці} & 6 & 0.860 & 0.860 & 0.860 & 0.016 \\
\hline
\specialcell{3cm}{По всій \\ вибірці} & 7 & 0.833 & 0.700 & 0.761 & 0.016 \\
\hline
\specialcell{3cm}{По всій \\ вибірці} & 8 & 0.913 & 0.840 & 0.875 & 0.016 \\
\hline
\specialcell{3cm}{По всій \\ вибірці} & 9 & 0.880 & 0.880 & 0.880 & 0.016 \\
\hline
\specialcell{3cm}{По всій \\ вибірці} & 10 & 0.913 & 0.840 & 0.875 & 0.016 \\
\hline
\specialcell{3cm}{По всій \\ вибірці} & 11 & 0.741 & 0.800 & 0.769 & 0.016 \\
\hline
\specialcell{3cm}{По всій \\ вибірці} & 12 & 1.000 & 0.980 & 0.990 & 0.016 \\
\hline
\specialcell{3cm}{По всій \\ вибірці} & 13 & 0.822 & 0.740 & 0.779 & 0.016 \\
\hline
\specialcell{3cm}{По всій \\ вибірці} & 14 & 0.792 & 0.760 & 0.776 & 0.016 \\
\hline
\specialcell{3cm}{По всій \\ вибірці} & 15 & 0.796 & 0.860 & 0.827 & 0.016 \\
\hline
\specialcell{3cm}{По всій \\ вибірці} & 16 & 0.870 & 0.800 & 0.833 & 0.016 \\
\hline
\specialcell{3cm}{По всій \\ вибірці} & 17 & 0.906 & 0.960 & 0.932 & 0.016 \\
\hline
\specialcell{3cm}{По всій \\ вибірці} & 18 & 0.885 & 0.920 & 0.902 & 0.016 \\
\hline
\specialcell{3cm}{По всій \\ вибірці} & 19 & 0.915 & 0.860 & 0.887 & 0.016 \\
\hline
\specialcell{3cm}{По всій \\ вибірці} & 20 & 0.980 & 0.960 & 0.970 & 0.016 \\
\hline
\specialcell{3cm}{По всій \\ вибірці} & 21 & 0.855 & 0.940 & 0.895 & 0.016 \\
\hline
\specialcell{3cm}{По всій \\ вибірці} & 22 & 0.868 & 0.920 & 0.893 & 0.016 \\
\hline
\specialcell{3cm}{По всій \\ вибірці} & 23 & 0.896 & 0.860 & 0.878 & 0.016 \\
\hline
\specialcell{3cm}{По всій \\ вибірці} & 24 & 0.894 & 0.840 & 0.866 & 0.016 \\
\hline
\specialcell{3cm}{По всій \\ вибірці} & 25 & 0.917 & 0.880 & 0.898 & 0.016 \\
\hline
\specialcell{3cm}{По всій \\ вибірці} & 26 & 0.837 & 0.820 & 0.828 & 0.016 \\
\hline
\specialcell{3cm}{По всій \\ вибірці} & 27 & 0.808 & 0.840 & 0.824 & 0.016 \\
\hline
\specialcell{3cm}{По всій \\ вибірці} & 28 & 0.885 & 0.920 & 0.902 & 0.016 \\
\hline
\specialcell{3cm}{По всій \\ вибірці} & 29 & 0.891 & 0.980 & 0.933 & 0.016 \\
\hline
\specialcell{3cm}{По всій \\ вибірці} & 30 & 0.907 & 0.980 & 0.942 & 0.016 \\
\hline
\specialcell{3cm}{По всій \\ вибірці} & 31 & 0.915 & 0.860 & 0.887 & 0.016 \\
\hline
\specialcell{3cm}{По всій \\ вибірці} & 32 & 0.904 & 0.940 & 0.922 & 0.016 \\
\hline
\specialcell{3cm}{По всій \\ вибірці} & 33 & 0.875 & 0.980 & 0.925 & 0.016 \\
\hline
\specialcell{3cm}{По всій \\ вибірці} & 34 & 0.918 & 0.900 & 0.909 & 0.016 \\
\hline
\specialcell{3cm}{По всій \\ вибірці} & 35 & 0.902 & 0.920 & 0.911 & 0.016 \\
\hline
\specialcell{3cm}{По всій \\ вибірці} & 36 & 1.000 & 0.900 & 0.947 & 0.016 \\
\hline
\specialcell{3cm}{По всій \\ вибірці} & 37 & 0.898 & 0.880 & 0.889 & 0.016 \\
\hline
\specialcell{3cm}{По всій \\ вибірці} & 38 & 0.840 & 0.840 & 0.840 & 0.016 \\
\hline
\specialcell{3cm}{По всій \\ вибірці} & 39 & 0.880 & 0.880 & 0.880 & 0.016 \\
\hline
\specialcell{3cm}{По всій \\ вибірці} & 40 & 0.878 & 0.860 & 0.869 & 0.016 \\
\hline
\specialcell{3cm}{По всій \\ вибірці} & 41 & 0.979 & 0.920 & 0.948 & 0.016 \\
\hline
\specialcell{3cm}{По всій \\ вибірці} & 42 & 0.957 & 0.900 & 0.928 & 0.016 \\
\hline
\specialcell{3cm}{По всій \\ вибірці} & 43 & 0.868 & 0.920 & 0.893 & 0.016 \\
\hline
\specialcell{3cm}{По всій \\ вибірці} & 44 & 0.632 & 0.720 & 0.673 & 0.016 \\
\hline
\specialcell{3cm}{По всій \\ вибірці} & 45 & 0.855 & 0.940 & 0.895 & 0.016 \\
\hline
\specialcell{3cm}{По всій \\ вибірці} & 46 & 0.920 & 0.920 & 0.920 & 0.016 \\
\hline
\specialcell{3cm}{По всій \\ вибірці} & 47 & 0.959 & 0.940 & 0.949 & 0.016 \\
\hline
\specialcell{3cm}{По всій \\ вибірці} & 48 & 0.907 & 0.980 & 0.942 & 0.016 \\
\hline
\specialcell{3cm}{По всій \\ вибірці} & 49 & 0.980 & 1.000 & 0.990 & 0.016 \\
\hline
\specialcell{3cm}{По всій \\ вибірці} & 50 & 0.941 & 0.960 & 0.950 & 0.016 \\
\hline
\specialcell{3cm}{По всій \\ вибірці} & 51 & 0.980 & 0.960 & 0.970 & 0.016 \\
\hline
\specialcell{3cm}{По всій \\ вибірці} & 52 & 1.000 & 0.980 & 0.990 & 0.016 \\
\hline
\specialcell{3cm}{По всій \\ вибірці} & 53 & 1.000 & 0.980 & 0.990 & 0.016 \\
\hline
\specialcell{3cm}{По всій \\ вибірці} & 54 & 1.000 & 1.000 & 1.000 & 0.016 \\
\hline
\specialcell{3cm}{По всій \\ вибірці} & 55 & 0.939 & 0.920 & 0.929 & 0.016 \\
\hline
\specialcell{3cm}{По всій \\ вибірці} & 56 & 0.923 & 0.960 & 0.941 & 0.016 \\
\hline
\specialcell{3cm}{По всій \\ вибірці} & 57 & 0.980 & 0.960 & 0.970 & 0.016 \\
\hline
\specialcell{3cm}{По всій \\ вибірці} & 58 & 0.889 & 0.960 & 0.923 & 0.016 \\
\hline
\specialcell{3cm}{По всій \\ вибірці} & 59 & 0.870 & 0.940 & 0.904 & 0.016 \\
\hline
\specialcell{3cm}{По всій \\ вибірці} & 60 & 0.959 & 0.940 & 0.949 & 0.016 \\
\hline
\specialcell{3cm}{По всій \\ вибірці} & 61 & 0.925 & 0.980 & 0.951 & 0.016 \\
\hline
\specialcell{3cm}{По всій \\ вибірці} & 62 & 0.759 & 0.820 & 0.788 & 0.016 \\
\hline
\specialcell{3cm}{По всій \\ вибірці} & 63 & 0.902 & 0.920 & 0.911 & 0.016 \\
\hline
\specialcell{3cm}{Тестовий \\ контекст} & 1 & 1.000 & 0.980 & 0.990 & 0.333 \\
\hline
\specialcell{3cm}{Тестовий \\ контекст} & 2 & 0.961 & 0.980 & 0.970 & 0.333 \\
\hline
\specialcell{3cm}{Тестовий \\ контекст} & 3 & 0.960 & 0.960 & 0.960 & 0.333 \\
\end{mytable*}

\chapter{Матриці помилок розпізнавання по реакціях}

\newcounter{contextnumber}

\begin{figure}
	\centering
	\subbottom[1 \label{img:confusion_matrix_data1_irs13_context_1}]{%
		\includegraphics[width=0.3\linewidth]{confusion_matrix_data1_irs13_context_1}}
	\forloop{contextnumber}{3}{\value{contextnumber} < 14}{
		\subbottom[\arabic{contextnumber} \label{img:confusion_matrix_data1_irs13_context_\arabic{contextnumber}}]{%
			\includegraphics[width=0.3\linewidth]{confusion_matrix_data1_irs13_context_\arabic{contextnumber}}}
	}
	
	\caption{Матриці помилок розпізнавання по реакціях для різних контекстів першого набору даних методом ІРС розміром 1--3 (частина 1)}
	\label{img:confusion_matrix_data1_irs13_1}
\end{figure}
\begin{figure}
	\centering
	\forloop{contextnumber}{14}{\value{contextnumber} < 20}{
		\subbottom[\arabic{contextnumber} \label{img:confusion_matrix_data1_irs13_context_\arabic{contextnumber}}]{%
			\includegraphics[width=0.3\linewidth]{confusion_matrix_data1_irs13_context_\arabic{contextnumber}}}
	}
	
	\caption{Матриці помилок розпізнавання по реакціях для різних контекстів першого набору даних методом ІРС розміром 1--3 (частина 2)}
	\label{img:confusion_matrix_data1_irs13_2}
\end{figure}
\begin{figure}
	\centering
	\includegraphics[width=\linewidth]{confusion_matrix_data1_cnn_context_20}
	\caption{Матриці помилок розпізнавання по реакціях по всій вибірці першого набору даних методом ІРС розміром 1--3}
	\label{img:confusion_matrix_data1_irs13_context_20}
\end{figure}

\begin{figure}
	\centering
	\subbottom[1 \label{img:confusion_matrix_data1_irs24_context_1}]{%
		\includegraphics[width=0.3\linewidth]{confusion_matrix_data1_irs24_context_1}}
	\forloop{contextnumber}{3}{\value{contextnumber} < 14}{
		\subbottom[\arabic{contextnumber} \label{img:confusion_matrix_data1_irs24_context_\arabic{contextnumber}}]{%
			\includegraphics[width=0.3\linewidth]{confusion_matrix_data1_irs24_context_\arabic{contextnumber}}}
	}
	
	\caption{Матриці помилок розпізнавання по реакціях для різних контекстів першого набору даних методом ІРС розміром 2--4 (частина 1)}
	\label{img:confusion_matrix_data1_irs24_1}
\end{figure}
\begin{figure}
	\centering
	\forloop{contextnumber}{14}{\value{contextnumber} < 20}{
		\subbottom[\arabic{contextnumber} \label{img:confusion_matrix_data1_irs24_context_\arabic{contextnumber}}]{%
			\includegraphics[width=0.3\linewidth]{confusion_matrix_data1_irs24_context_\arabic{contextnumber}}}
	}
	
	\caption{Матриці помилок розпізнавання по реакціях для різних контекстів першого набору даних методом ІРС розміром 2--4 (частина 2)}
	\label{img:confusion_matrix_data1_irs24_2}
\end{figure}
\begin{figure}
	\centering
	\includegraphics[width=\linewidth]{confusion_matrix_data1_cnn_context_20}
	\caption{Матриці помилок розпізнавання по реакціях по всій вибірці першого набору даних методом ІРС розміром 2--4}
	\label{img:confusion_matrix_data1_irs24_context_20}
\end{figure}

\begin{figure}
	\centering
	\subbottom[1 \label{img:confusion_matrix_data1_cnn_context_1}]{%
		\includegraphics[width=0.3\linewidth]{confusion_matrix_data1_cnn_context_1}}
	\forloop{contextnumber}{3}{\value{contextnumber} < 14}{
		\subbottom[\arabic{contextnumber} \label{img:confusion_matrix_data1_cnn_context_\arabic{contextnumber}}]{%
			\includegraphics[width=0.3\linewidth]{confusion_matrix_data1_cnn_context_\arabic{contextnumber}}}
	}
	
	\caption{Матриці помилок розпізнавання по реакціях для різних контекстів першого набору даних методом ЗНМ (частина 1)}
	\label{img:confusion_matrix_data1_cnn_1}
\end{figure}
\begin{figure}
	\centering
	\forloop{contextnumber}{14}{\value{contextnumber} < 20}{
		\subbottom[\arabic{contextnumber} \label{img:confusion_matrix_data1_cnn_context_\arabic{contextnumber}}]{%
			\includegraphics[width=0.3\linewidth]{confusion_matrix_data1_cnn_context_\arabic{contextnumber}}}
	}
	
	\caption{Матриці помилок розпізнавання по реакціях для різних контекстів першого набору даних методом ЗНМ (частина 2)}
	\label{img:confusion_matrix_data1_cnn_2}
\end{figure}
\begin{figure}
	\centering
	\includegraphics[width=\linewidth]{confusion_matrix_data1_cnn_context_20}
	\caption{Матриці помилок розпізнавання по реакціях по всій вибірці першого набору даних методом ЗНМ}
	\label{img:confusion_matrix_data1_cnn_context_20}
\end{figure}

\begin{figure}
	\centering
	\subbottom[1 \label{img:confusion_matrix_data3_irs13_context_1}]{%
		\includegraphics[width=0.3\linewidth]{confusion_matrix_data3_irs13_context_1}}
	\forloop{contextnumber}{3}{\value{contextnumber} < 14}{
		\subbottom[\arabic{contextnumber} \label{img:confusion_matrix_data3_irs13_context_\arabic{contextnumber}}]{%
			\includegraphics[width=0.3\linewidth]{confusion_matrix_data3_irs13_context_\arabic{contextnumber}}}
	}
	
	\caption{Матриці помилок розпізнавання по реакціях для різних контекстів третього набору даних методом ІРС розміром 1--3 (частина 1)}
	\label{img:confusion_matrix_data3_irs13_1}
\end{figure}
\begin{figure}
	\centering
	\forloop{contextnumber}{14}{\value{contextnumber} < 20}{
		\subbottom[\arabic{contextnumber} \label{img:confusion_matrix_data3_irs13_context_\arabic{contextnumber}}]{%
			\includegraphics[width=0.3\linewidth]{confusion_matrix_data3_irs13_context_\arabic{contextnumber}}}
	}
	
	\caption{Матриці помилок розпізнавання по реакціях для різних контекстів третього набору даних методом ІРС розміром 1--3 (частина 2)}
	\label{img:confusion_matrix_data3_irs13_2}
\end{figure}
\begin{figure}
	\centering
	\includegraphics[width=\linewidth]{confusion_matrix_data3_cnn_context_20}
	\caption{Матриці помилок розпізнавання по реакціях по всій вибірці третього набору даних методом ІРС розміром 1--3}
	\label{img:confusion_matrix_data3_irs13_context_20}
\end{figure}

\begin{figure}
	\centering
	\subbottom[1 \label{img:confusion_matrix_data3_irs24_context_1}]{%
		\includegraphics[width=0.3\linewidth]{confusion_matrix_data3_irs24_context_1}}
	\forloop{contextnumber}{3}{\value{contextnumber} < 14}{
		\subbottom[\arabic{contextnumber} \label{img:confusion_matrix_data3_irs24_context_\arabic{contextnumber}}]{%
			\includegraphics[width=0.3\linewidth]{confusion_matrix_data3_irs24_context_\arabic{contextnumber}}}
	}
	
	\caption{Матриці помилок розпізнавання по реакціях для різних контекстів третього набору даних методом ІРС розміром 2--4 (частина 1)}
	\label{img:confusion_matrix_data3_irs24_1}
\end{figure}
\begin{figure}
	\centering
	\forloop{contextnumber}{14}{\value{contextnumber} < 20}{
		\subbottom[\arabic{contextnumber} \label{img:confusion_matrix_data3_irs24_context_\arabic{contextnumber}}]{%
			\includegraphics[width=0.3\linewidth]{confusion_matrix_data3_irs24_context_\arabic{contextnumber}}}
	}
	
	\caption{Матриці помилок розпізнавання по реакціях для різних контекстів третього набору даних методом ІРС розміром 2--4 (частина 2)}
	\label{img:confusion_matrix_data3_irs24_2}
\end{figure}
\begin{figure}
	\centering
	\includegraphics[width=\linewidth]{confusion_matrix_data3_cnn_context_20}
	\caption{Матриці помилок розпізнавання по реакціях по всій вибірці третього набору даних методом ІРС розміром 2--4}
	\label{img:confusion_matrix_data3_irs24_context_20}
\end{figure}

\begin{figure}
	\centering
	\subbottom[1 \label{img:confusion_matrix_data3_cnn_context_1}]{%
		\includegraphics[width=0.3\linewidth]{confusion_matrix_data3_cnn_context_1}}
	\forloop{contextnumber}{3}{\value{contextnumber} < 14}{
		\subbottom[\arabic{contextnumber} \label{img:confusion_matrix_data3_cnn_context_\arabic{contextnumber}}]{%
			\includegraphics[width=0.3\linewidth]{confusion_matrix_data3_cnn_context_\arabic{contextnumber}}}
	}
	
	\caption{Матриці помилок розпізнавання по реакціях для різних контекстів третього набору даних методом ЗНМ (частина 1)}
	\label{img:confusion_matrix_data3_cnn_1}
\end{figure}
\begin{figure}
	\centering
	\forloop{contextnumber}{14}{\value{contextnumber} < 20}{
		\subbottom[\arabic{contextnumber} \label{img:confusion_matrix_data3_cnn_context_\arabic{contextnumber}}]{%
			\includegraphics[width=0.3\linewidth]{confusion_matrix_data3_cnn_context_\arabic{contextnumber}}}
	}
	
	\caption{Матриці помилок розпізнавання по реакціях для різних контекстів третього набору даних методом ЗНМ (частина 2)}
	\label{img:confusion_matrix_data3_cnn_2}
\end{figure}
\begin{figure}
	\centering
	\includegraphics[width=\linewidth]{confusion_matrix_data3_cnn_context_20}
	\caption{Матриці помилок розпізнавання по реакціях по всій вибірці третього набору даних методом ЗНМ}
	\label{img:confusion_matrix_data3_cnn_context_20}
\end{figure}

\chapter{Довідка про впровадження} \label{AppendixС}

\begin{center}
	\includegraphics[width=\linewidth]{190104170842}
\end{center}

\chapter{Список публікацій здобувача за темою дисертації та відомості про апробацію результатів дисертації}

\printbibliography[heading=pubsubgroup, keyword=biblioauthorvak, notcategory=biblioauthoreng, section=1,title=\vakbibtitle]%
\printbibliography[heading=pubsubgroup, keyword=biblioauthorvak, category=biblioauthoreng, section=1,title=\vakindexedbibtitle]%
\printbibliography[heading=pubsubgroup, notkeyword=biblioauthorvak, category=biblioauthoreng, section=1,title=\engbibtitle]%
\printbibliography[heading=pubsubgroup, keyword=biblioauthornotvak, notcategory=biblioauthoreng, section=1,title=\notvakbibtitle]%
\printbibliography[heading=pubsubgroup, keyword=biblioauthorconf, section=1,title=\confbibtitle]%

\begin{center}
	\textit{Основні положення дисертаційної роботи були апробовані у насткпних міжнародних науково-практичних конференціях}
\end{center}

\begin{itemize}
	\item XVII Мiжнародна науково-технiчна конференцiя «Системний аналiз та iнформацiйнi технологiї» (м. Київ, 2015 р.)
	\item XIII Міжнародна конференція «Управління проектами у розвитку суспільства», тема: «Проекти в умовах глобальних загроз, ризиків и викликів» (м. Київ, 2016 р.)
	\item ІІІ Міжнародна науково-практична конференція «Інформаційні технології та взаємодії» (м. Київ, 2016 р.)
	\item 16th EAGE International Conference on Geoinformatics - Theoretical and Applied Aspects (м. Київ, 2017 р.)
	\item ІV Міжнародна науково-практична конференція «Інформаційні технології та взаємодії» (м. Київ, 2017 р.)
\end{itemize}